% Blueprint content for the AbsorptionCutoff formalization.
% Nodes carry \lean{...} (Lean declaration), \leanok (formalized), and
% \uses{...} (dependencies) so leanblueprint can build the dependency graph
% and checkdecls can verify the declaration names against the Lean project.

\chapter{Setup: rounding, the squared-radius mean map, and cutoff}

This chapter collects the state-independent scaffolding of the paper (§1--§2):
the nearest-grid rounding map, the squared-radius mean map $V_A$ and its
fixed-point structure, and the total-variation cutoff / mixing-time definitions.
It corresponds to the first formalized batch.

\section{Nearest-grid rounding}

\begin{definition}[Unit-grid rounding]
  \label{def:Q1}
  \lean{AbsorptionCutoff.Q₁}
  \leanok
  The scalar nearest-integer rounding $Q_1:\mathbb{R}\to\mathbb{Z}$, with ties
  broken toward the grid point of smaller absolute value (round half toward
  zero).
\end{definition}

\begin{lemma}[Zero bin]
  \label{lem:Q1-zero}
  \lean{AbsorptionCutoff.Q₁_zero_iff}
  \leanok
  \uses{def:Q1}
  $Q_1 u = 0 \iff |u| \le \tfrac12$.
\end{lemma}
\begin{proof}
  \leanok
  \uses{def:Q1}
  Split on the sign of $u$ and use $\lceil r\rceil = 0 \iff r\in(-1,0]$.
\end{proof}

\begin{lemma}[Rounding error]
  \label{lem:Q1-err}
  \lean{AbsorptionCutoff.Q₁_sub_le}
  \leanok
  \uses{def:Q1}
  $|Q_1 u - u| \le \tfrac12$.
\end{lemma}
\begin{proof}
  \leanok
  \uses{def:Q1}
  From $r \le \lceil r\rceil < r+1$ with $r = |u| - \tfrac12$.
\end{proof}

\begin{definition}[Grid and vector rounding]
  \label{def:Qrho}
  \lean{AbsorptionCutoff.gridRound, AbsorptionCutoff.Qρ}
  \leanok
  \uses{def:Q1}
  $\mathrm{gridRound}_\varrho(u) = \varrho\, Q_1(u/\varrho)$ and the
  coordinatewise vector rounding $Q_\varrho(x)_i = \mathrm{gridRound}_\varrho(x_i)$
  onto $(\varrho\mathbb{Z})^N$.
\end{definition}

\begin{lemma}[Vector zero bin]
  \label{lem:Qrho-zero}
  \lean{AbsorptionCutoff.Qρ_eq_zero_iff}
  \leanok
  \uses{def:Qrho}
  For $\varrho>0$, $Q_\varrho x = 0 \iff |x_i|\le \varrho/2$ for all $i$.
\end{lemma}
\begin{proof}\leanok\uses{lem:Q1-zero}\end{proof}

\section{The squared-radius mean map}

\begin{definition}[Squared-radius mean map]
  \label{def:V}
  \lean{AbsorptionCutoff.V}
  \leanok
  $V_A(q) = \mathbb{E}\!\left[\tanh^2\!\big(A\sqrt{q}\,G\big)\right]$, where
  $G\sim\mathcal N(0,1)$.
\end{definition}

\begin{lemma}[Elementary values]
  \label{lem:V-basic}
  \lean{AbsorptionCutoff.V_zero, AbsorptionCutoff.V_nonneg}
  \leanok
  \uses{def:V}
  $V_A(0) = 0$ and $V_A(q) \ge 0$.
\end{lemma}
\begin{proof}\leanok\uses{def:V}\end{proof}

\begin{theorem}[Unique positive fixed point]
  \label{thm:fixed-point}
  \lean{AbsorptionCutoff.exists_unique_positive_fixed}
  \leanok
  \uses{lem:no-two-fixed}
  A continuous, strictly concave self-map $f$ of $[0,1]$ with $f(0)=0$,
  $f(1)<1$, that exceeds the diagonal at some interior point, has a unique fixed
  point in $(0,1)$. This is the state-independent core of the paper's
  Lemma on fixed points and linearization of $V_A$.
\end{theorem}
\begin{proof}
  \leanok
  \uses{lem:no-two-fixed}
  Existence: the intermediate value theorem applied to $g=f-\mathrm{id}$ on
  $[q_0,1]$, where $g(q_0)>0>g(1)$. Uniqueness: Lemma~\ref{lem:no-two-fixed}.
\end{proof}

\begin{lemma}[At most one positive fixed point]
  \label{lem:no-two-fixed}
  \lean{AbsorptionCutoff.no_two_fixed}
  \leanok
  A strictly concave $f$ on $[0,1]$ with $f(0)=0$ has no two distinct positive
  fixed points.
\end{lemma}
\begin{proof}
  \leanok
  Writing the smaller fixed point $a$ as a strict convex combination of $0$ and
  the larger $b$, strict concavity gives $f(a) > a$, contradicting $f(a)=a$.
\end{proof}

\begin{lemma}[$V_A < 1$]
  \label{lem:V-lt-one}
  \lean{AbsorptionCutoff.V_lt_one}
  \leanok
  \uses{def:V}
  $V_A(q) < 1$ for all $q$; in particular $V_A(1)<1$. Since $\tanh^2<1$ pointwise,
  $1-V_A(q)=\int(1-\tanh^2)\,\mathrm d\mu>0$.
\end{lemma}
\begin{proof}\leanok\uses{def:V}\end{proof}

\begin{lemma}[$V_A$ continuous]
  \label{lem:V-continuous}
  \lean{AbsorptionCutoff.V_continuous}
  \leanok
  \uses{def:V}
  $V_A$ is continuous (continuity under the integral sign, dominated by $1$).
\end{lemma}
\begin{proof}\leanok\uses{def:V}\end{proof}

\begin{theorem}[Fixed point of $V_A$, conditional]
  \label{thm:V-fixed}
  \lean{AbsorptionCutoff.V_exists_unique_fixed_of_structural}
  \leanok
  \uses{thm:fixed-point, def:V, lem:V-basic, lem:V-lt-one, lem:V-continuous}
  Conditional only on strict concavity of $V_A$, the map has a unique positive fixed
  point; continuity and the boundary value $V_A(1)<1$ are supplied by
  Lemmas~\ref{lem:V-continuous} and~\ref{lem:V-lt-one}. The unconditional $A>1$
  statement follows after the strict-concavity and near-zero slope estimates below.
\end{theorem}
\begin{proof}\leanok\uses{thm:fixed-point, lem:V-lt-one, lem:V-continuous}\end{proof}

\begin{lemma}[$\tanh x / x$ strictly decreasing]
  \label{lem:tanh-div}
  \lean{AbsorptionCutoff.tanh_div_self_strictAntiOn}
  \leanok
  The map $x\mapsto \tanh x/x$ is strictly decreasing on $(0,\infty)$.
\end{lemma}
\begin{proof}[Proof outline]
  \leanok
  Its derivative is $((1-\tanh^2 x)x-\tanh x)/x^2$, which is negative because
  $x<\sinh(2x)/2$.
\end{proof}

\begin{lemma}[Ratio monotonicity of $V_A$]
  \label{lem:V-ratio}
  \lean{AbsorptionCutoff.V_ratio_strictAntiOn}
  \leanok
  \uses{lem:tanh-div, def:V}
  For $A\neq0$, $q\mapsto V_A(q)/q$ is strictly decreasing on $(0,1)$: integrate the
  pointwise decrease of $\tanh^2(A\sqrt q\,g)/q$, strict off the Gaussian-null set
  $\{g=0\}$.
\end{lemma}
\begin{proof}\leanok\uses{lem:tanh-div}\end{proof}

\begin{lemma}[$V_A$ strictly concave]
  \label{lem:V-concave}
  \lean{AbsorptionCutoff.V_strictConcaveOn}
  \leanok
  \uses{def:V, lem:tanh-div}
  For $A\neq0$, $V_A$ is strictly concave on $[0,1]$ (paper's
  \emph{lem:gaussian-mean-field-concavity}). Pointwise $q\mapsto\tanh^2(A g\sqrt q)$
  is strictly concave (its derivative $a^2(\tanh x/x)\sech^2 x$, $x=a\sqrt q$, is
  strictly decreasing); integrate, strict off the Gaussian-null set $\{g=0\}$.
\end{lemma}
\begin{proof}\leanok\uses{lem:tanh-div}\end{proof}

\begin{lemma}[$V_A$ strictly increasing]
  \label{lem:V-mono}
  \lean{AbsorptionCutoff.V_strictMonoOn}
  \leanok
  \uses{def:V}
  For $A\neq0$, $V_A$ is strictly increasing on $[0,1]$: pointwise
  $q\mapsto\tanh^2(A g\sqrt q)$ is increasing, strict off $\{g=0\}$; integrate.
\end{lemma}
\begin{proof}\leanok\uses{def:V}\end{proof}

\begin{theorem}[Unique fixed point of $V_A$]
  \label{thm:V-fixed-uncond}
  \lean{AbsorptionCutoff.V_exists_unique_fixed}
  \leanok
  \uses{lem:V-concave, lem:V-continuous, lem:V-lt-one}
  For $A\neq0$, if $V_A$ exceeds the diagonal at some $q_0\in(0,1)$ (true for $A>1$),
  it has a unique positive fixed point. This is the paper's
  \emph{lem:gaussian-mean-field-concavity}, proved the paper's way via strict concavity
  of $V_A$ (continuity and $V_A(1)<1$ are also discharged).
\end{theorem}
\begin{proof}\leanok\uses{lem:V-concave}\end{proof}

\begin{lemma}[$V_A$ differentiable under the integral]
  \label{lem:V-deriv}
  \lean{AbsorptionCutoff.hasDerivAt_V}
  \leanok
  \uses{def:V}
  For $A\neq0$ and $q>0$, $V_A$ is differentiable at $q$ with
  $V_A'(q)=\int f_{Ag}'(q)\,d\mu$, obtained by differentiating under the integral sign.
  The derivative integrand $f_{Ag}'(q)=(Ag)^2(\tanh x/x)\sech^2 x$ ($x=A\sqrt q\,g$) is
  dominated by $g\mapsto(Ag)^2$, integrable since the standard Gaussian has a finite
  second moment.
\end{lemma}
\begin{proof}\leanok\end{proof}

\begin{lemma}[Fixed-point multiplier $\mu_A\in(0,1)$]
  \label{lem:V-multiplier}
  \lean{AbsorptionCutoff.V_multiplier_mem_Ioo}
  \leanok
  \uses{lem:V-deriv, lem:V-concave}
  At any positive fixed point $q_\ast\in(0,1)$ of $V_A$ ($A\neq0$), the multiplier
  $\mu_A:=V_A'(q_\ast)$ satisfies $0<\mu_A<1$ (paper \emph{eq:gaussian-fixed-point}):
  positivity from $V_A'(q_\ast)=\int f_{Ag}'(q_\ast)>0$ (integrand positive off
  $\{g=0\}$); the upper bound from strict concavity, the graph lying below the tangent
  at $q_\ast$, so $V_A'(q_\ast)<V_A(q_\ast)/q_\ast=1$.
\end{lemma}
\begin{proof}\leanok\uses{lem:V-deriv, lem:V-concave}\end{proof}

\begin{lemma}[Near-zero slope $V_A'(0+)=A^2$]
  \label{lem:V-deriv-zero}
  \lean{AbsorptionCutoff.V_ratio_tendsto}
  \leanok
  \uses{def:V}
  For $A\neq0$, $V_A(q)/q\to A^2$ as $q\to0^+$ (paper
  \emph{eq:gaussian-mean-field-near-zero}). Dominated convergence of
  $g\mapsto\tanh^2(A\sqrt q\,g)/q\to(Ag)^2$ (bound $\tanh^2 x\le x^2$, dominating
  function $A^2g^2$), using the Gaussian second moment $\int g^2=1$.
\end{lemma}
\begin{proof}\leanok\end{proof}

\begin{theorem}[Unconditional fixed point for $A>1$]
  \label{thm:V-fixed-gt-one}
  \lean{AbsorptionCutoff.V_exists_unique_fixed_of_one_lt}
  \leanok
  \uses{lem:V-deriv-zero, thm:V-fixed-uncond}
  For $A>1$, $V_A$ has a unique positive fixed point $q_\ast\in(0,1)$, with no
  exceed-diagonal hypothesis: since $V_A'(0+)=A^2>1$, $V_A$ exceeds the diagonal near
  $0$ (\texttt{V\_exceeds\_diagonal\_of\_one\_lt}), discharging the last hypothesis of
  \cref{thm:V-fixed-uncond}. This is the full $A>1$ case of
  \emph{lem:gaussian-mean-field-concavity}.
\end{theorem}
\begin{proof}\leanok\uses{lem:V-deriv-zero}\end{proof}

\begin{theorem}[Only fixed point $0$ for $A\le1$]
  \label{thm:V-fixed-le-one}
  \lean{AbsorptionCutoff.V_fixed_eq_zero_of_le_one}
  \leanok
  \uses{lem:V-ratio, lem:V-deriv-zero, lem:V-lt-one}
  For $0<A\le1$, the only fixed point of $V_A$ in $[0,1]$ is $0$ (the $A\le1$ case of
  \emph{lem:gaussian-mean-field-concavity}). Since $q\mapsto V_A(q)/q$ is strictly
  decreasing with right-limit $A^2$ at $0^+$, it is $<A^2\le1$ throughout $(0,1)$
  (\texttt{V\_ratio\_lt\_sq}), so $V_A(q)<q$ there (\texttt{V\_lt\_self\_of\_le\_one});
  with $V_A(1)<1$ this excludes any positive fixed point.
\end{theorem}
\begin{proof}\leanok\uses{lem:V-ratio, lem:V-deriv-zero}\end{proof}

\begin{lemma}[Orbit sign relative to the fixed point]
  \label{lem:V-orbit-sign}
  \lean{AbsorptionCutoff.V_gt_self_of_lt_fixed, AbsorptionCutoff.V_lt_self_of_gt_fixed}
  \leanok
  \uses{lem:V-ratio}
  Let $q_\ast\in(0,1)$ be a fixed point of $V_A$ ($A\ne0$). Since $q\mapsto V_A(q)/q$ is
  strictly decreasing and equals $1$ at $q_\ast$, one has $V_A(q)>q$ for $q\in(0,q_\ast)$
  and $V_A(q)<q$ for $q\in(q_\ast,1)$.
\end{lemma}
\begin{proof}\leanok\uses{lem:V-ratio}\end{proof}

\begin{theorem}[Deterministic orbit convergence $V_A^t(q)\to q_\ast$]
  \label{thm:V-orbit-convergence}
  \lean{AbsorptionCutoff.V_orbit_tendsto, AbsorptionCutoff.exists_fixed_and_orbit_tendsto,
    AbsorptionCutoff.V_orbit_tendsto_Ioc,
    AbsorptionCutoff.summable_abs_V_orbit_sub_fixed,
    AbsorptionCutoff.koenigsOrbitFactor, AbsorptionCutoff.V_iterate_ne_fixed,
    AbsorptionCutoff.summable_abs_koenigsOrbitFactor_sub_one,
    AbsorptionCutoff.koenigsOrbitFactor_pos, AbsorptionCutoff.koenigsOrbitProduct,
    AbsorptionCutoff.multipliable_koenigsOrbitFactor_and_product_ne_zero,
    AbsorptionCutoff.mul_prod_range_koenigsOrbitFactor_eq_normalized,
    AbsorptionCutoff.mul_koenigsOrbitProduct_ne_zero,
    AbsorptionCutoff.tendsto_normalized_V_orbit_sub_fixed,
    AbsorptionCutoff.koenigsCoefficient,
    AbsorptionCutoff.koenigsCoefficient_ne_zero,
    AbsorptionCutoff.V_orbit_sub_fixed_sub_koenigsCoefficient_mul_pow_isLittleO,
    AbsorptionCutoff.exists_uniform_eventual_geometric_V_orbit_bound,
    AbsorptionCutoff.exists_uniform_summable_koenigsOrbitFactor_tail_bound,
    AbsorptionCutoff.exists_hasProdUniformlyOn_koenigsOrbitFactor_tail,
    AbsorptionCutoff.continuousOn_koenigsOrbitFactor,
    AbsorptionCutoff.exists_tendstoUniformlyOn_koenigsCoefficient,
    AbsorptionCutoff.continuousOn_koenigsCoefficient,
    AbsorptionCutoff.exists_tendstoUniformlyOn_shifted_normalized_V_orbit_sub_fixed,
    AbsorptionCutoff.exists_tendstoUniformlyOn_normalized_V_orbit_sub_fixed,
    AbsorptionCutoff.exists_tendstoUniformlyOn_koenigs_remainder_ratio,
    AbsorptionCutoff.exists_continuous_local_uniform_koenigs_asymptotic}
  \leanok
  \uses{lem:V-orbit-sign, lem:V-mono, lem:V-continuous, lem:V-lt-one, thm:V-fixed-gt-one}
  Assume $A>1$, and let $q_\ast$ be the unique positive fixed point of $V_A$ and
  $\mu_A=V_A'(q_\ast)\in(0,1)$. For every $q\in(0,1]$,
  \[
    V_A^n(q)\longrightarrow q_\ast,
    \qquad
    \sum_{n=0}^\infty |V_A^n(q)-q_\ast|<\infty.
  \]
  If $q\ne q_\ast$, there is a nonzero coefficient $C(q)$ such that
  \[
    V_A^n(q)-q_\ast-C(q)\mu_A^n=o(\mu_A^n)
  \]
  as $n\to\infty$. The coefficient $C$ is continuous and nonzero locally on
  $(0,1]$ away from $q_\ast$, and the normalized remainder
  \[
    \frac{V_A^n(q)-q_\ast-C(q)\mu_A^n}{\mu_A^n}
  \]
  tends to zero locally uniformly there.
\end{theorem}
\begin{proof}[Proof outline]
  \leanok
  \uses{lem:V-orbit-sign, lem:V-mono, lem:V-continuous, thm:V-fixed-gt-one}
  The orbit-sign lemma makes every orbit monotone and bounded. Its limit is a
  fixed point by continuity, and hence equals $q_\ast$. After finite entrance
  into a neighborhood on which $V_A$ is a strict contraction, the orbit error
  has a summable geometric bound.

  For $q\ne q_\ast$, strict monotonicity keeps the orbit on one side of the
  fixed point. The positive secant factors are therefore well-defined, their
  deviations from one are summable, and their infinite product is nonzero.
  The finite-product identity gives the coefficient $C(q)$ and the pointwise
  asymptotic. On a neighborhood of any $q\ne q_\ast$, a common entrance time
  and summable majorant give uniform convergence of the product. This proves
  continuity of $C$ and the locally uniform remainder estimate.
\end{proof}

\begin{lemma}[Deterministic entrance into the stable interval]
  \label{lem:gaussian-deterministic-entrance}
  \lean{AbsorptionCutoff.exists_local_V_contraction_with_deriv,
    AbsorptionCutoff.abs_V_sub_fixed_le_abs_sub_fixed,
    AbsorptionCutoff.antitone_abs_V_iterate_sub_fixed,
    AbsorptionCutoff.exists_uniform_V_iterate_abs_sub_fixed_lt,
    AbsorptionCutoff.exists_uniform_V_iterate_stable_margin,
    AbsorptionCutoff.exists_uniform_stable_V_orbit_interval_Icc,
    AbsorptionCutoff.exists_pos_uniform_V_tracking_tolerance,
    AbsorptionCutoff.exists_uniform_markovPathMeasure_stable_entrance_exp_bound,
    AbsorptionCutoff.exists_uniform_Kchain_pow_stable_entrance_exp_bound,
    AbsorptionCutoff.exists_uniform_stable_V_orbit_interval,
    AbsorptionCutoff.exists_eventually_stable_V_orbit_interval}
  \leanok
  \uses{thm:V-orbit-convergence, lem:V-continuous, lem:V-deriv}
  A compact interval centered at the positive fixed point can be chosen
  strictly inside $(0,1)$ so that one factor $\kappa<1$ bounds both
  $|V_A'|$ and the map contraction throughout it. For every macroscopic
  initial radius, one neighborhood enters a smaller concentric interval at a
  common time and remains a positive distance from the outer boundary, with a
  uniform geometric bound thereafter. Consequently, if $q_N\to q_0$ and
  eventually $q_N\in(0,1]$, one dimension threshold $N_0$ makes these
  conclusions hold for every $N\ge N_0$, exactly as in the paper's
  \texttt{eq:gaussian-deterministic-entrance} and
  \texttt{eq:gaussian-stable-basin-derivative}.
\end{lemma}
\begin{proof}\leanok\uses{thm:V-orbit-convergence, lem:V-continuous, lem:V-deriv}\end{proof}

\chapter{Common scalar estimates}

\begin{lemma}[Negative moments from Laplace envelopes]
  \label{lem:laplace-envelope-negative-moment}
  \lean{AbsorptionCutoff.integrable_neg_rpow_and_integral_le_of_laplace_le,
    AbsorptionCutoff.exp_neg_add_one_mul_self_rpow_le_Gamma_add_one,
    AbsorptionCutoff.rpow_div_Gamma_add_one_le_exp_mul_div_rpow}
  \leanok
  Let $X$ be a measurable random variable with $X>0$ almost surely and let
  $p>0$. If its Laplace transform is bounded by $B(t)$ for $t>0$ and
  $t^{p-1}B(t)$ is integrable there, then $X^{-p}$ is integrable and
  \[
    \mathbb E[X^{-p}]
    \le \Gamma(p)^{-1}\int_0^\infty t^{p-1}B(t)\,dt.
  \]
  The proof applies Fubini to the Mellin--Laplace representation. Almost-sure
  strict positivity is essential because Lean's real-power convention assigns
  value zero to $0^{-p}$. Integrating the Gamma density only over
  $[p,p+1]$ also gives the elementary lower bound
  $e^{-(p+1)}p^p\leq\Gamma(p+1)$, which is sufficient for exponential
  dimension estimates and avoids invoking Stirling's formula. In normalized
  form this gives
  $n^p/\Gamma(p+1)\leq e^{p+1}(n/p)^p$.
\end{lemma}
\begin{proof}\leanok\end{proof}

\begin{lemma}[Hölder interpolation on a measurable set]
  \label{lem:set-fractional-moment-holder}
  \lean{AbsorptionCutoff.integrableOn_rpow_and_setIntegral_rpow_le}
  \leanok
  Let $f\ge0$ almost everywhere be integrable on a finite measure space.
  For a measurable set $S$ and $0<\theta<1$,
  $f^\theta$ is integrable on $S$ and
  \[
    \int_S f^\theta\,d\mu
    \le
    \left(\int f\,d\mu\right)^\theta
    \mu(S)^{1-\theta}.
  \]
  This is the reusable Hölder step for controlling lower negative moments on
  the truncated-Cramér bad event.
\end{lemma}
\begin{proof}\leanok\end{proof}

\begin{lemma}[Hölder interpolation for negative moments]
  \label{lem:set-negative-moment-holder}
  \lean{AbsorptionCutoff.integrableOn_neg_rpow_and_setIntegral_neg_rpow_le}
  \leanok
  Let $X>0$ almost everywhere on a finite measure space, let $N>0$, and let
  $0<\gamma<\alpha$. If $X^{-\alpha N}$ is integrable, then for every
  measurable set $S$, the function $X^{-\gamma N}$ is integrable on $S$ and
  \[
    \int_S X^{-\gamma N}\,d\mu
    \le
    \left(\int X^{-\alpha N}\,d\mu\right)^{\gamma/\alpha}
    \mu(S)^{1-\gamma/\alpha}.
  \]
\end{lemma}
\begin{proof}\leanok\end{proof}

\begin{lemma}[Preserving a strict product gap]
  \label{lem:strict-product-gap}
  \lean{AbsorptionCutoff.exists_eps_delta_of_one_lt_mul}
  \leanok
  If $a>0$ and $a\mu>1$, then there are
  $\varepsilon\in(0,1)$ and $\delta>0$ such that
  \[
    (1-\varepsilon)a(\mu-\delta)>1.
  \]
  An explicit two-stage construction first spends half of the additive
  mean gap on $\delta$, then half of the remaining multiplicative gap on
  $\varepsilon$.
\end{lemma}
\begin{proof}\leanok\end{proof}

\begin{lemma}[Killed-moment recursion bound]
  \label{lem:geom-recursion}
  \lean{AbsorptionCutoff.geom_recursion_bound,
    AbsorptionCutoff.geom_recursion_bound_contraction,
    AbsorptionCutoff.geom_recursion_bound_contraction_of_lt,
    AbsorptionCutoff.geom_recursion_bound_contraction_pow_of_lt,
    AbsorptionCutoff.eventually_nat_mul_exp_neg_le_inv_nat,
    AbsorptionCutoff.eventually_two_nat_horizon_mul_exp_neg_le_mul_inv_nat}
  \leanok
  Iterating $m_{t+1} \le a_t m_t + b_t$ (with $a_u \ge 0$) gives
  \[ m_t \le \Big(\textstyle\prod_{u<t} a_u\Big) m_0
      + \sum_{s<t}\Big(\textstyle\prod_{s<u<t} a_u\Big) b_s. \]
  This is the deterministic engine of the paper's killed second-moment bound;
  the stochastic wrapper is deferred. For a constant
  coefficient $0\le a<1$ and increment $B$, direct induction also gives the
  uniform fixed-point bound
  $m_t\le\max\{m_0,B/(1-a)\}$.
\end{lemma}
\begin{proof}\leanok\end{proof}

\begin{lemma}[Uniform killed-moment bound]
  \label{lem:geom-recursion-const}
  \lean{AbsorptionCutoff.geom_recursion_bound_const}
  \leanok
  \uses{lem:geom-recursion}
  With $a_u \le \alpha$ ($\alpha\ge1$) and $b_s \le \beta$, the bound simplifies to
  $m_t \le \alpha^t m_0 + t\,\alpha^t\beta$ — the form used with the amplification
  bound $K_N$.
\end{lemma}
\begin{proof}\leanok\uses{lem:geom-recursion}\end{proof}

\begin{lemma}[Backward weights for the stopped-moment telescoping]
  \label{lem:back-weight}
  \lean{AbsorptionCutoff.backWeight, AbsorptionCutoff.backWeight_self, AbsorptionCutoff.one_le_backWeight,
    AbsorptionCutoff.backWeight_succ_le, AbsorptionCutoff.sq_mul_backWeight_succ_le}
  \leanok
  For the stopped-moment estimate, define the backward weights
  $W_t=\prod_{u=t}^{T-1}(1\vee\alpha_u^2)$. Each factor is $\ge1$, so $W_T=1$,
  $1\le W_t$, the weights are nonincreasing ($W_{t+1}\le W_t$), and they dominate the
  forward multiplier: $\alpha_t^2W_{t+1}\le W_t$. These feed the telescoping
  $\Ee[W_{t+1}\overline{\mathcal R}_{t+1}^2]-\Ee[W_t\overline{\mathcal R}_t^2]\le CK_N^2/N$;
  the stochastic wrapper (filtration, stopping time, conditional expectation) is deferred.
\end{lemma}
\begin{proof}\leanok\end{proof}

\begin{lemma}[Killed second-moment estimate (stochastic wrapper)]
  \label{lem:killed-moment}
  \lean{AbsorptionCutoff.killed_step_pointwise, AbsorptionCutoff.condExp_add_sq_le,
    AbsorptionCutoff.killed_condExp_step, AbsorptionCutoff.killedMoment,
    AbsorptionCutoff.killed_moment_recursion, AbsorptionCutoff.killed_moment_closed_form}
  \leanok
  \uses{lem:geom-recursion}
  On a filtered probability space, let $\mathcal R$ be adapted, let $A_t$ be
  $\mathcal F_t$-measurable with $|A_t|\le\alpha_t$, and let $\eta$ be
  martingale-difference noise satisfying
  ($\Ee[\eta_{t+1}\mid\mathcal F_t]=0$, $\Ee[\eta_{t+1}^2\mid\mathcal F_t]\le C\sigma_t^2/N$),
  and let $\tau$ be a stopping time. Then
  \[
    m_{t+1}\leq\alpha_t^2m_t+\frac{C\sigma_t^2}{N},
    \qquad
    m_t=\Ee[\mathcal R_t^2\mathbf 1_{\{\tau>t\}}].
  \]
  The recursion bound of \Cref{lem:geom-recursion} gives the corresponding
  closed-form killed-moment estimate.
\end{lemma}
\begin{proof}[Proof outline]
  \leanok
  \uses{lem:geom-recursion}
  Pointwise,
  $\mathbf{1}_{\tau>t+1}\mathcal R_{t+1}^2\le
  \mathbf{1}_{\tau>t}(A_t\mathcal R_t+\eta_{t+1})^2$. Taking
  $\Ee[\,\cdot\mid\mathcal F_t]$ (the cross term drops since $\Ee[\eta\mid\mathcal F_t]=0$)
  gives $\Ee[\mathbf{1}_{\tau>t+1}\mathcal R_{t+1}^2\mid\mathcal F_t]\le
  \mathbf{1}_{\tau>t}(\alpha_t^2\mathcal R_t^2+C\sigma_t^2/N)$.
  Integrate and apply \Cref{lem:geom-recursion}.
\end{proof}

\begin{lemma}[Weighted stopped one-step estimate]
  \label{lem:stopped-condexp-step}
  \lean{AbsorptionCutoff.stoppedValue, AbsorptionCutoff.stoppedValue_eq_stoppedProcess_coe,
    AbsorptionCutoff.stoppedValue_succ_decomposition,
    AbsorptionCutoff.stoppedValue_succ_sq_decomposition,
    AbsorptionCutoff.condExp_indicator_add_sq_le, AbsorptionCutoff.stopped_condExp_step,
    AbsorptionCutoff.integral_le_of_condExp_add_const, AbsorptionCutoff.stopped_moment_recursion,
    AbsorptionCutoff.sum_recursion_bound, AbsorptionCutoff.stopped_moment_telescope,
    AbsorptionCutoff.sum_weighted_noise_le, AbsorptionCutoff.stopped_moment_uniform_bound,
    AbsorptionCutoff.trackingAmplification, AbsorptionCutoff.tracking_product_le,
    AbsorptionCutoff.tracking_noise_product_le,
    AbsorptionCutoff.backWeight_zero_le_trackingAmplification_sq,
    AbsorptionCutoff.backWeight_mul_sigma_sq_le_trackingAmplification_sq,
    AbsorptionCutoff.stopped_moment_backWeight_endpoints,
    AbsorptionCutoff.stopped_moment_tracking_bound,
    AbsorptionCutoff.exit_indicator_sq_le_stopped,
    AbsorptionCutoff.exit_prob_mul_sq_le_stopped_moment,
    AbsorptionCutoff.exit_prob_le_stopped_moment_div,
    AbsorptionCutoff.exit_prob_le_stopped_moment_div_of_le,
    AbsorptionCutoff.orbit_exit_prob_bound}
  \leanok
  \uses{lem:back-weight}
  For backward weights satisfying $0\le W_{t+1}\le W_t$ and
  $\alpha_t^2W_{t+1}\le W_t$, the stopped recursion obeys
  \[
    \Ee[W_{t+1}\overline{\mathcal R}_{t+1}^2\mid\mathcal F_t]
    \le W_t\overline{\mathcal R}_t^2+W_{t+1}C\sigma_t^2/N.
  \]
  The $\{\tau\le t\}$ branch is already stopped and is
  $\mathcal F_t$-measurable; on $\{t<\tau\}$ the martingale-difference cross term
  vanishes after conditioning. Summing the scalar recursion telescopes. The finite-horizon
  amplification scalar $K$ bounds both $W_0^{1/2}$ and
  $\sigma_tW_{t+1}^{1/2}$, yielding
  $\Ee\overline{\mathcal R}_T^2\le CK^2(\Ee\mathcal R_0^2+T/N)$. If $\tau$ is the
  first exit from $[-\delta,\delta]$, then
  $\delta^2\mathbf 1_{\{\tau\le T\}}\le\overline{\mathcal R}_T^2$ pointwise;
  integration gives
  $\Pp(\tau\le T)\le CK^2\delta^{-2}(\Ee\mathcal R_0^2+T/N)$.
\end{lemma}
\begin{proof}\leanok\uses{lem:back-weight}\end{proof}

\section{Total-variation cutoff and mixing time}

\begin{definition}[Total-variation distance]
  \label{def:tvdist}
  \lean{AbsorptionCutoff.tvDist}
  \leanok
  $\|\mu-\nu\|_{\mathrm{TV}} = \sup_{B}|\mu(B)-\nu(B)|$ over measurable $B$.
\end{definition}

\begin{definition}[Distance to equilibrium]
  \label{def:dseq}
  \lean{AbsorptionCutoff.dSeq}
  \leanok
  \uses{def:tvdist}
  $d(t) = \|\kappa^t(x,\cdot) - \pi\|_{\mathrm{TV}}$, using the Mathlib kernel
  power $\kappa^t$.
\end{definition}

\begin{definition}[Admissible cutoff center and window]
  \label{def:cutoff-window}
  \lean{AbsorptionCutoff.IsCutoffWindow}
  \leanok
  A cutoff center and window are admissible when $t_n\to\infty$, the window
  $a_n$ is eventually positive, and $a_n=o(t_n)$. Eventual positivity is the
  sequence formulation of the paper's $a_n>0$: changing finitely many indices
  does not affect any cutoff limit.
\end{definition}

\begin{definition}[Windowed cutoff limits]
  \label{def:cutoff-limits}
  \lean{AbsorptionCutoff.HasCutoffLimits}
  \leanok
  \uses{def:dseq}
  The two early and late iterated total-variation limits obtained by evaluating
  at $t_n-ca_n$ and $t_n+ca_n$, respectively, and then sending $c\to\infty$.
\end{definition}

\begin{definition}[Total-variation cutoff]
  \label{def:cutoff}
  \lean{AbsorptionCutoff.HasCutoff}
  \leanok
  \uses{def:cutoff-window, def:cutoff-limits}
  The conjunction of admissibility of $(t_n,a_n)$ and the two windowed cutoff
  limits. Thus this is the full cutoff definition used in the paper, rather than
  only its two limiting conclusions.
\end{definition}

\begin{definition}[Mixing time]
  \label{def:mixing}
  \lean{AbsorptionCutoff.mixingTime}
  \leanok
  \uses{def:dseq}
  $t_{\mathrm{mix}}(\varepsilon) = \inf\{t : d(t)\le\varepsilon\}$, with value
  $\top$ if no such $t$ exists.
\end{definition}

\section{The absorption identity}

\begin{definition}[Absorbing state]
  \label{def:absorbing}
  \lean{AbsorptionCutoff.IsAbsorbing}
  \leanok
  A point $a$ is absorbing for a kernel $\kappa$ if $\kappa(a,\cdot)=\delta_a$; the
  rounded chain's origin is absorbing.
\end{definition}

\begin{lemma}[Absorption time and canonical Markov trajectory]
  \label{lem:absorption-time}
  \lean{AbsorptionCutoff.absorptionTime, AbsorptionCutoff.Adapted.isStoppingTime_absorptionTime,
    AbsorptionCutoff.absorptionTime_gt_iff, AbsorptionCutoff.zero_persists, AbsorptionCutoff.zero_persists',
    AbsorptionCutoff.absorptionTime_gt_iff_value_ne_zero,
    AbsorptionCutoff.measure_absorptionTime_gt_eq,
    AbsorptionCutoff.measure_absorptionTime_gt_eq_of_ae,
    AbsorptionCutoff.measure_absorptionTime_gt_eq_of_map,
    AbsorptionCutoff.measure_absorptionTime_gt_eq_kernel_pow,
    AbsorptionCutoff.measure_absorptionTime_gt_eq_of_ae_map,
    AbsorptionCutoff.measure_absorptionTime_gt_eq_of_ae_kernel_pow,
    AbsorptionCutoff.markovHistoryKernel, AbsorptionCutoff.markovPathMeasure,
    AbsorptionCutoff.markovPathMeasure_map_zero, AbsorptionCutoff.markovPathMeasure_map_eval_succ,
    AbsorptionCutoff.markovPathMeasure_map_eval, AbsorptionCutoff.markovPathMeasure_dirac_map_eval,
    AbsorptionCutoff.markovPathMeasure_ae_absorbing}
  \leanok
  \uses{def:absorbing}
  The first hitting time of zero is a stopping time for an adapted process. If zero is
  pathwise absorbing, then $\{\tau>t\}=\{X_t\ne0\}$, so any path law whose time-$t$
  marginal is $\kappa^t(x,\cdot)$ satisfies
  $\Pp(\tau>t)=\kappa^t(x,\{0\}^c)$. The canonical homogeneous path law is built from
  Mathlib's Ionescu--Tulcea \texttt{trajMeasure}; its time-$t$ coordinate marginal is
  identified with the kernel power, $(\text{\texttt{markovPathMeasure}}\,\delta_x\,\kappa)
  \circ (\cdot_t)^{-1} = \kappa^t(x,\cdot)$, by induction from the time-zero marginal and
  the one-step recursion obtained from the Ionescu--Tulcea joint law. On the raw canonical
  path space pathwise absorption fails, so the survival bridge is used in its
  almost-everywhere form: from $\kappa(0,\cdot)=\delta_0$ the coordinate process is
  absorbing $\Pp$-a.e. (\texttt{markovPathMeasure\_ae\_absorbing}, via the Ionescu--Tulcea
  joint law), whence $\{\tau>t\}=\{X_t\ne0\}$ a.e.\ and the same marginal-law identity
  $\Pp(\tau>t)=\kappa^t(x,\{0\}^c)$ holds.
\end{lemma}
\begin{proof}\leanok\uses{def:absorbing}\end{proof}

\begin{lemma}[TV to a point mass is the survival mass]
  \label{lem:tv-dirac}
  \lean{AbsorptionCutoff.tvDist_dirac}
  \leanok
  \uses{def:tvdist}
  For a probability measure $\mu$, $\|\mu-\delta_a\|_{\mathrm{TV}} = \mu(\{a\}^c)$.
  Applied to $\mu = P^t(x,\cdot)$ with $a$ the absorbing origin, this is the core of
  the paper's identity $\|P^t(x,\cdot)-\delta_0\|_{\mathrm{TV}} = \mathbb{P}(\tau > t)$.
\end{lemma}
\begin{proof}
  \leanok
  \uses{def:tvdist}
  Both bounds are attained at $B=\{a\}^c$: every measurable $B$ gives
  $|\mu(B)-\delta_a(B)| \le \mu(\{a\}^c)$, with equality there.
\end{proof}

\begin{lemma}[Absorption mass is monotone]
  \label{lem:absorb-mono}
  \lean{AbsorptionCutoff.absorb_mass_mono}
  \leanok
  \uses{def:absorbing, def:dseq}
  For an absorbing state, $\kappa^t(x,\{a\})$ is nondecreasing in $t$; hence the
  survival mass and $d(t)$ are nonincreasing (paper `rem:intro-profile-mixing`).
\end{lemma}
\begin{proof}
  \leanok
  \uses{def:absorbing}
  Since $\kappa(y,\{a\})\geq\mathbf 1_{\{a\}}(y)$,
  \[
    \kappa^{t+1}(x,\{a\})
      =\int \kappa(y,\{a\})\,\kappa^t(x,\mathrm dy)
      \geq\int \mathbf 1_{\{a\}}(y)\,\kappa^t(x,\mathrm dy)
      =\kappa^t(x,\{a\}).
  \]
\end{proof}

\begin{lemma}[Kernel form of the absorption identity]
  \label{lem:tv-pow-dirac}
  \lean{AbsorptionCutoff.tvDist_pow_dirac}
  \leanok
  \uses{lem:tv-dirac, def:dseq}
  For a Markov kernel and an absorbing origin,
  $\|\kappa^t(x,\cdot)-\delta_a\|_{\mathrm{TV}} = \kappa^t(x,\{a\}^c) = \mathbb{P}(\tau>t)$.
\end{lemma}
\begin{proof}\leanok\uses{lem:tv-dirac}\end{proof}

\begin{lemma}[Distance to equilibrium is nonincreasing]
  \label{lem:dseq-antitone}
  \lean{AbsorptionCutoff.dSeq_dirac_antitone}
  \leanok
  \uses{lem:tv-pow-dirac, lem:absorb-mono}
  For an absorbing state, $t\mapsto d(t)$ is nonincreasing, so the mixing times of
  Definition~\ref{def:mixing} are well behaved.
\end{lemma}
\begin{proof}\leanok\uses{lem:absorb-mono}\end{proof}

\begin{lemma}[TV metric facts for probability measures]
  \label{lem:tv-metric}
  \lean{AbsorptionCutoff.tvDist_comm, AbsorptionCutoff.tvDist_le_one, AbsorptionCutoff.tvDist_bddAbove}
  \leanok
  \uses{def:tvdist}
  $\|\cdot\|_{\mathrm{TV}}$ is symmetric, and for probability measures it is bounded
  by $1$ (so the defining supremum is $\mathrm{BddAbove}$).
\end{lemma}
\begin{proof}\leanok\uses{def:tvdist}\end{proof}

\chapter{The squared-radius Markov chains}

This chapter realizes the paper's §2 squared-radius chains as Mathlib Markov kernels
and connects the rounded chain to the abstract absorption identity of the previous
chapter. The step maps act on the squared radius $q\in[0,1]$; the driving randomness
is a standard Gaussian vector.

\begin{definition}[Gaussian driving noise]
  \label{def:gaussianVec}
  \lean{AbsorptionCutoff.gaussianVec}
  \leanok
  $\mathcal G_N = \mathcal N(0,1)^{\otimes N}$ on $\mathbb R^N$, the product of $N$
  standard Gaussians (a probability measure).
\end{definition}

\begin{definition}[Unrounded step map]
  \label{def:Fmap}
  \lean{AbsorptionCutoff.Fmap}
  \leanok
  \uses{def:gaussianVec}
  $F_{A,N}(q,g) = \tfrac1N\sum_{i<N}\tanh^2\!\big(A\sqrt q\,g_i\big)$, the empirical
  squared radius after one Gaussian step.
\end{definition}

\begin{lemma}[Step map takes values in $[0,1)$]
  \label{lem:Fmap-range}
  \lean{AbsorptionCutoff.Fmap_nonneg, AbsorptionCutoff.Fmap_lt_one}
  \leanok
  \uses{def:Fmap}
  $0 \le F_{A,N}(q,g)$, and $F_{A,N}(q,g) < 1$ for $N>0$ (each $\tanh^2 < 1$).
\end{lemma}
\begin{proof}\leanok\uses{def:Fmap}\end{proof}

\begin{definition}[Unrounded squared-radius kernel]
  \label{def:Kchain}
  \lean{AbsorptionCutoff.Kchain}
  \leanok
  \uses{def:Fmap}
  $K_{A,N}(q,\cdot) = (F_{A,N}(q,\cdot))_\#\mathcal G_N$, the Markov kernel pushing the
  Gaussian noise through the step map. (\texttt{measurable\_Fmap} supplies
  measurability; the kernel is Markov since $\mathcal G_N$ is a probability measure.)
\end{definition}

\begin{lemma}[Mean of the unrounded step is $V_A$]
  \label{lem:integral-Fmap}
  \lean{AbsorptionCutoff.integral_Fmap}
  \leanok
  \uses{def:Fmap, def:V}
  $\displaystyle\int F_{A,N}(q,g)\,\mathrm d\mathcal G_N(g) = V_A(q)$ for $N\neq0$
  (the paper's \texttt{eq:gaussian-q-moments}). Each coordinate marginal of
  $\mathcal G_N$ is $\mathcal N(0,1)$, so the $N$ identical terms average to
  $\mathbb E[\tanh^2(A\sqrt q\,G)] = V_A(q)$.
\end{lemma}
\begin{proof}\leanok\uses{def:V}\end{proof}

\begin{lemma}[Unrounded one-step Hoeffding bound]
  \label{lem:Fmap-hoeffding}
  \lean{AbsorptionCutoff.FmapSubgaussianParameter,
    AbsorptionCutoff.coe_FmapSubgaussianParameter,
    AbsorptionCutoff.hasSubgaussianMGF_Fmap_sub_V,
    AbsorptionCutoff.measureReal_abs_Fmap_sub_V_gt_le,
    AbsorptionCutoff.Kchain_measureReal_abs_sub_V_gt_le,
    AbsorptionCutoff.markovPathMeasure_ae_eval_succ_mem_Kchain_Icc,
    AbsorptionCutoff.markovPathMeasure_dirac_ae_eval_mem_Kchain_Icc,
    AbsorptionCutoff.markovPathMeasure_dirac_ae_forall_le_mem_Kchain_Icc,
    AbsorptionCutoff.markovPathMeasure_measureReal_abs_next_sub_V_gt_le,
    AbsorptionCutoff.finiteHorizonKchainStepDeviationEvent,
    AbsorptionCutoff.markovPathMeasure_measureReal_finiteHorizonKchainStepDeviationEvent_le,
    AbsorptionCutoff.exists_pos_markovPathMeasure_measureReal_abs_sub_iterate_V_ge_le}
  \leanok
  \uses{def:Fmap, lem:integral-Fmap}
  The centered empirical update is sub-Gaussian with exact proxy
  $1/(4N)$. Consequently, uniformly in the current radius $q$,
  \[
    \mathbb P\!\left(\left|F_{A,N}(q,G)-V_A(q)\right|>\varepsilon\right)
    \leq 2e^{-2N\varepsilon^2}.
  \]
\end{lemma}
\begin{proof}\leanok\uses{lem:integral-Fmap}\end{proof}

\begin{definition}[Rounded step map]
  \label{def:Hmap}
  \lean{AbsorptionCutoff.Hmap}
  \leanok
  \uses{def:Q1, def:gaussianVec}
  The rounded analogue of $F_{A,N}$ is
  \[
    H_{A,\varrho,N}(h,g)
      =\frac1N\sum_{i<N}
        \left(\varrho^{-1}Q_1\!\left(
          \varrho^{-1}\tanh(\varrho A\sqrt h\,g_i)
        \right)\right)^2.
  \]
\end{definition}

\begin{lemma}[Rounded step vanishes at the origin]
  \label{lem:Hmap-zero}
  \lean{AbsorptionCutoff.Hmap_zero, AbsorptionCutoff.Hmap_nonneg}
  \leanok
  \uses{def:Hmap, lem:Q1-zero}
  $H_{A,\varrho,N}(0,g)=0$ (since $\tanh 0 = 0$ lands in the zero bin, $Q_1(0)=0$),
  and $H_{A,\varrho,N}\ge0$.
\end{lemma}
\begin{proof}\leanok\uses{lem:Q1-zero}\end{proof}

\begin{definition}[Rounded squared-radius kernel]
  \label{def:Hkernel}
  \lean{AbsorptionCutoff.Hkernel}
  \leanok
  \uses{def:Hmap}
  $H_{A,\varrho,N}(h,\cdot) = (H_{A,\varrho,N}(h,\cdot))_\#\mathcal G_N$, the rounded
  Markov kernel (\texttt{measurable\_Hmap}, via measurability of $Q_1$).
\end{definition}

\begin{lemma}[The origin is absorbing for the rounded chain]
  \label{lem:Hkernel-absorbing}
  \lean{AbsorptionCutoff.isAbsorbing_Hkernel}
  \leanok
  \uses{def:Hkernel, def:absorbing, lem:Hmap-zero}
  $H_{A,\varrho,N}(0,\cdot) = \delta_0$: the rounded step map is identically $0$ at
  $h=0$, so its pushforward of any probability measure is $\delta_0$.
\end{lemma}
\begin{proof}\leanok\uses{lem:Hmap-zero}\end{proof}

\begin{lemma}[Concrete rounded absorption identities and decay]
  \label{lem:Hkernel-tv}
  \lean{AbsorptionCutoff.tvDist_Hkernel_pow_dirac, AbsorptionCutoff.dSeq_Hkernel_dirac_antitone,
    AbsorptionCutoff.measure_roundedAbsorptionTime_gt_eq, AbsorptionCutoff.tvDist_Hkernel_pow_eq_survival,
    AbsorptionCutoff.isAbsorbing_roundedPkernel,
    AbsorptionCutoff.measure_roundedVectorAbsorptionTime_gt_eq,
    AbsorptionCutoff.tvDist_roundedPkernel_pow_dirac,
    AbsorptionCutoff.tvDist_roundedPkernel_pow_eq_survival,
    AbsorptionCutoff.tvDist_roundedPkernel_pow_eq_Hkernel_pow}
  \leanok
  \uses{lem:Hkernel-absorbing, lem:tv-pow-dirac, lem:dseq-antitone, lem:absorption-time}
  For the rounded chain,
  $\|H^t_{A,\varrho,N}(x,\cdot)-\delta_0\|_{\mathrm{TV}} = H^t_{A,\varrho,N}(x,\{0\}^c)
  = \mathbb P(\tau_\varrho > t)$, and $t\mapsto d(t)$ is nonincreasing — the paper's
  \texttt{eq:tv-absorption} instantiated for the concrete chain. Running the chain on its
  canonical Ionescu--Tulcea path law (\Cref{lem:absorption-time}) makes $\tau_\varrho$ the
  first hitting time of $0$ by the coordinate process, and the a.e.\ survival bridge gives
  $\mathbb P(\tau_\varrho > t)=H^t_{A,\varrho,N}(x,\{0\}^c)$ literally on that space, whence
  $\|H^t_{A,\varrho,N}(x,\cdot)-\delta_0\|_{\mathrm{TV}} = \mathbb P(\tau_\varrho > t)$.
  The same identities hold for the rounded vector kernel $P_{\varrho,A,N}$ on its
  canonical path space.  Moreover, its normalized rounded squared-radius pushforward is
  exactly the $H_{A,\varrho,N}$ law, so the vector and scalar TV distances from their
  absorbing origins agree at every time.
\end{lemma}
\begin{proof}\leanok\uses{lem:Hkernel-absorbing, lem:tv-pow-dirac, lem:absorption-time}\end{proof}

\begin{lemma}[Lattice-subcritical contraction]
  \label{lem:subcritical-lattice-regimes}
  \lean{AbsorptionCutoff.latticeProfile, AbsorptionCutoff.latticeMap, AbsorptionCutoff.roundedMeanMap,
    AbsorptionCutoff.latticeProfile_ratio_tendsto_zero,
    AbsorptionCutoff.latticeProfile_ratio_tendsto_one,
    AbsorptionCutoff.latticeThreshold, AbsorptionCutoff.latticeThreshold_pos,
    AbsorptionCutoff.latticeThreshold_lt_one, AbsorptionCutoff.abs_Q₁_mono,
    AbsorptionCutoff.exists_latticeMap_le_mul, AbsorptionCutoff.roundedMeanMap_le_latticeMap,
    AbsorptionCutoff.exists_roundedMeanMap_le_mul}
  \leanok
  \uses{lem:Q1-zero, lem:Q1-err}
  Let $m(\alpha)=\mathbb E[Q_1(\alpha G)^2]$,
  $L_A(h)=m(A\sqrt h)$, and
  $A_{\mathrm{lat}}^2=\inf_{\alpha>0}\alpha^2/m(\alpha)$.
  The endpoint limits $m(\alpha)/\alpha^2\to0$ as $\alpha\downarrow0$ and
  $m(\alpha)/\alpha^2\to1$ as $\alpha\to\infty$, together with the paper's
  positive lower bound and strict test at $\alpha=1/2$, give
  $A_{\mathrm{lat}}\in(0,1)$. If $0<A<A_{\mathrm{lat}}$, compact-middle
  maximization and the endpoint limits yield
  $L_A(h)\leq\theta_Ah$ for all $h>0$, for some $\theta_A\in(0,1)$.
  Finally, $|Q_1|$ is monotone in the magnitude of its argument and
  $|\tanh x|\leq|x|$, hence
  $V_{A,\varrho}(h)\leq L_A(h)\leq\theta_Ah$ for every $h\geq0$ and $\varrho>0$.
\end{lemma}
\begin{proof}\leanok\uses{lem:Q1-zero, lem:Q1-err}\end{proof}

\begin{lemma}[Exact fixed-precision layer formulas]
  \label{lem:subcritical-layer-formulas-exact}
  \lean{AbsorptionCutoff.roundedLayerThreshold, AbsorptionCutoff.roundedLayerIndices,
    AbsorptionCutoff.roundedMeanMap_eq_sum_upperTail,
    AbsorptionCutoff.hasDerivAt_gaussianUpperTail,
    AbsorptionCutoff.hasDerivAt_two_mul_gaussianUpperTail_div_sqrt,
    AbsorptionCutoff.hasDerivAt_roundedMeanMap}
  \leanok
  \uses{lem:subcritical-lattice-regimes}
  For $h>0$, the rounded mean map is the finite sum of Gaussian upper tails over the
  admissible thresholds $b_{\varrho,k}$, and its derivative is obtained term by term:
  \[
    V'_{A,\varrho}(h)=\sum_{k\in\mathcal K_\varrho}(2k+1)
    \frac{b_{\varrho,k}}{Ah^{3/2}}
    \phi\!\left(\frac{b_{\varrho,k}}{A\sqrt h}\right).
  \]
\end{lemma}
\begin{proof}\leanok\uses{lem:subcritical-lattice-regimes}\end{proof}

\begin{lemma}[Fixed-precision map estimates]
  \label{lem:subcritical-fixed-precision-map-estimates}
  \lean{AbsorptionCutoff.exists_roundedMeanMap_compact_bounds,
    AbsorptionCutoff.gaussianUpperTail_le_exp_neg_sq_div_two,
    AbsorptionCutoff.roundedMeanMap_le_weight_mul_exp,
    AbsorptionCutoff.gaussianPDFReal_layer_le_common_exp,
    AbsorptionCutoff.roundedMeanMapDerivative_le_weight_mul_rpow_exp,
    AbsorptionCutoff.exists_roundedMeanMap_small_radius_bounds}
  \leanok
  \uses{lem:subcritical-layer-formulas-exact}
  On every compact interval $[r,R]\subset(0,\infty)$, the rounded mean map has a
  positive lower bound and its derivative has a finite upper bound.  Moreover, for
  every $r\in(0,1)$ there are $c,C>0$ such that, for $0<h\leq r$,
  \[
    V_{A,\varrho}(h)\leq C e^{-c/h},
    \qquad
    V'_{A,\varrho}(h)\leq C h^{-2}e^{-c/h}.
  \]
\end{lemma}
\begin{proof}\leanok\uses{lem:subcritical-layer-formulas-exact}\end{proof}

\begin{lemma}[Normalized amplification along rounded deterministic orbits]
  \label{lem:subcritical-normalized-lattice-amplification}
  \lean{AbsorptionCutoff.roundedOrbit, AbsorptionCutoff.roundedLocalDerivativeSup,
    AbsorptionCutoff.roundedBeta, AbsorptionCutoff.roundedAlpha,
    AbsorptionCutoff.roundedAmplificationProduct,
    AbsorptionCutoff.roundedAmplificationCombination,
    AbsorptionCutoff.exists_eventually_roundedAmplificationCombination_le_exp_sq,
    AbsorptionCutoff.trackingAmplification_le_of_combination_le,
    AbsorptionCutoff.exists_eventually_uniform_subcritical_tail_length_le_mul_log_log,
    AbsorptionCutoff.exists_eventually_roundedAlpha_le_mul_log,
    AbsorptionCutoff.exists_eventually_uniform_amplificationProduct_le_exp_sq,
    AbsorptionCutoff.exists_eventually_uniform_roundedAmplificationCombination_le_exp_sq,
    AbsorptionCutoff.exists_eventually_uniform_trackingAmplification_le_exp_sq,
    AbsorptionCutoff.tendsto_inv_nat_mul_exp_log_log_sq_zero,
    AbsorptionCutoff.tendsto_succ_nat_div_mul_exp_log_log_sq_zero}
  \leanok
  \uses{lem:subcritical-lattice-regimes,
    lem:subcritical-fixed-precision-map-estimates}
  For every lattice-subcritical rounded orbit starting in a fixed compact
  interval, the product-plus-radius-ratio expression in
  \texttt{eq:subcritical-global-amplification} is bounded, uniformly over all
  nonempty time intervals, by
  $\exp\{C(\log\log N)^2\}$ for all sufficiently large $N$.
\end{lemma}
\begin{proof}
  \leanok
  \uses{lem:subcritical-lattice-regimes,
    lem:subcritical-fixed-precision-map-estimates}
\end{proof}

\begin{lemma}[Fixed-precision coordinate moment bound]
  \label{lem:subcritical-fixed-moment-bound}
  \lean{AbsorptionCutoff.roundedCoordinateObservable,
    AbsorptionCutoff.integral_roundedCoordinateObservable,
    AbsorptionCutoff.measureReal_roundedCoordinateObservable_ne_zero_le,
    AbsorptionCutoff.integral_pow_roundedCoordinateObservable_le_exp,
    AbsorptionCutoff.roundedSupportExponent_eq,
    AbsorptionCutoff.exists_pos_exp_neg_div_le_mul_pow,
    AbsorptionCutoff.exists_pos_integral_pow_roundedCoordinateObservable_le_mul_pow}
  \leanok
  \uses{lem:subcritical-fixed-precision-map-estimates}
  For every integer $p\geq1$, there is $C_p>0$ such that the one-coordinate
  rounded squared-radius observable satisfies
  \[
    \mathbb E\!\left[U_{A,\varrho}(h,G)^p\right]\leq C_p h^p,
    \qquad h\geq0.
  \]
  Indeed, nonzero rounded output forces $|G|$ beyond the first admissible layer,
  whose Gaussian probability is exponentially small in $1/h$.  The resulting
  exponential moment bound dominates $h^p$ near zero; boundedness of the rounded
  observable handles $h\geq1$.
\end{lemma}
\begin{proof}
  \leanok
  \uses{lem:subcritical-fixed-precision-map-estimates}
\end{proof}

\begin{lemma}[One-step rounded radius variance]
  \label{lem:subcritical-fixed-noise-variance}
  \lean{AbsorptionCutoff.integral_Hmap_eq_roundedMeanMap,
    AbsorptionCutoff.integral_sq_Hmap_sub_roundedMeanMap_le,
    AbsorptionCutoff.exists_pos_integral_sq_Hmap_sub_roundedMeanMap_le}
  \leanok
  \uses{lem:subcritical-fixed-moment-bound}
  Under the standard Gaussian product law, the rounded radius step has mean
  $V_{A,\varrho}(h)$ and satisfies
  \[
    \mathbb E\!\left[
      \bigl(H_{A,\varrho,N}(h,G)-V_{A,\varrho}(h)\bigr)^2
    \right]
    \leq \frac{C h^2}{N}.
  \]
  This follows from independence of the $N$ coordinate observables, the
  finite-product variance identity, and the $p=2$ coordinate moment bound.
\end{lemma}
\begin{proof}
  \leanok
  \uses{lem:subcritical-fixed-moment-bound}
\end{proof}

\begin{lemma}[Canonical conditional rounded-noise bound]
  \label{lem:subcritical-fixed-conditional-noise}
  \lean{AbsorptionCutoff.condDistrib_markovPathMeasure_eval_succ,
    AbsorptionCutoff.condExp_markovPathMeasure_prefix_eval_succ,
    AbsorptionCutoff.condExp_markovPathMeasure_prefix_eval_succ_piLE,
    AbsorptionCutoff.continuous_roundedMeanMap,
    AbsorptionCutoff.measurable_Hmap_right, AbsorptionCutoff.integrable_Hmap,
    AbsorptionCutoff.Hkernel_apply,
    AbsorptionCutoff.integral_Hkernel_sub_roundedMeanMap_eq_zero,
    AbsorptionCutoff.Hkernel_apply_Icc_compl,
    AbsorptionCutoff.markovPathMeasure_ae_eval_succ_mem_Icc,
    AbsorptionCutoff.integral_sq_Hkernel_sub_roundedMeanMap,
    AbsorptionCutoff.integrable_next_sub_roundedMeanMap,
    AbsorptionCutoff.condExp_next_sub_roundedMeanMap_eq_zero,
    AbsorptionCutoff.integrable_sq_next_sub_roundedMeanMap,
    AbsorptionCutoff.condExp_sq_next_sub_roundedMeanMap_eq,
    AbsorptionCutoff.exists_pos_condExp_sq_next_sub_roundedMeanMap_le,
    AbsorptionCutoff.exists_pos_condExp_sq_next_sub_roundedMeanMap_le_of_tracking}
  \leanok
  \uses{lem:subcritical-fixed-noise-variance}
  On the canonical path space for the rounded radius kernel, conditioning on
  the coordinate-prefix filtration gives
  \[
    \mathbb E\!\left[
      \bigl(\mathscr H_{t+1}^{(N)}
        -V_{A,\varrho}(\mathscr H_t^{(N)})\bigr)^2
      \,\middle|\,\mathcal F_t
    \right]
    \leq
    \frac{C}{N}\bigl(\mathscr H_t^{(N)}\bigr)^2
    \quad\text{a.s.}
  \]
  The Ionescu--Tulcea conditional-distribution identity supplies the conditional
  kernel, compact support discharges integrability, and the Gaussian pushforward
  formula reduces the conditional second moment to
  \Cref{lem:subcritical-fixed-noise-variance}.  On the tracking event
  $|\mathscr H_t^{(N)}-\bar h_t|\leq
  \delta(\bar h_t+a_N)$, the same bound and deterministic localization give
  \[
    \mathbb E[\xi_{t+1}^2\mid\mathcal F_t]
    \leq \frac{C(1+\delta)^2}{N}(\bar h_t+a_N)^2
    \quad\text{a.s.}
  \]
\end{lemma}
\begin{proof}
  \leanok
  \uses{lem:subcritical-fixed-noise-variance}
\end{proof}

\begin{lemma}[Canonical rounded-radius tracking recursion]
  \label{lem:subcritical-radius-tracking-recursion}
  \lean{AbsorptionCutoff.radiusTrackingError, AbsorptionCutoff.radiusNoise,
    AbsorptionCutoff.normalizedRadiusError,
    AbsorptionCutoff.radiusDifferenceQuotient,
    AbsorptionCutoff.normalizedRadiusNoise,
    AbsorptionCutoff.radiusNormalizedMultiplier,
    AbsorptionCutoff.radiusTrackingGood, AbsorptionCutoff.radiusTrackingExitTime,
    AbsorptionCutoff.radiusTrackingExitTime_le_sentinel,
    AbsorptionCutoff.lt_radiusTrackingExitTime_iff,
    AbsorptionCutoff.radiusTrackingGood_iff_abs_normalizedRadiusError_le,
    AbsorptionCutoff.radiusTrackingExitTime_le_iff_exists_abs_normalizedRadiusError_gt,
    AbsorptionCutoff.radiusTrackingGood_of_lt_exit,
    AbsorptionCutoff.delta_lt_abs_normalizedRadiusError_at_exit,
    AbsorptionCutoff.eval_mem_roundedOrbitWindow_of_lt_exit,
    AbsorptionCutoff.deriv_roundedMeanMap_nonneg_of_nonneg,
    AbsorptionCutoff.exists_deriv_roundedMeanMap_le_on_Icc_zero,
    AbsorptionCutoff.bddAbove_deriv_roundedMeanMap_image_roundedOrbitWindow,
    AbsorptionCutoff.uIcc_subset_roundedOrbitWindow_of_mem,
    AbsorptionCutoff.exists_deriv_eq_secant_roundedMeanMap,
    AbsorptionCutoff.exists_mem_roundedOrbitWindow_deriv_eq_radiusDifferenceQuotient,
    AbsorptionCutoff.radiusDifferenceQuotient_le_roundedLocalDerivativeSup_of_bddAbove,
    AbsorptionCutoff.radiusDifferenceQuotient_le_roundedLocalDerivativeSup,
    AbsorptionCutoff.abs_radiusNormalizedMultiplier_le_roundedAlpha_of_mem,
    AbsorptionCutoff.abs_radiusNormalizedMultiplier_le_roundedAlpha_of_lt_exit,
    AbsorptionCutoff.measurable_piLE_eval,
    AbsorptionCutoff.measurable_radiusTrackingError,
    AbsorptionCutoff.measurable_normalizedRadiusError,
    AbsorptionCutoff.measurable_radiusDifferenceQuotient,
    AbsorptionCutoff.measurable_radiusNormalizedMultiplier,
    AbsorptionCutoff.measurableSet_radiusTrackingGood,
    AbsorptionCutoff.measurableSet_lt_radiusTrackingExitTime,
    AbsorptionCutoff.measurableSet_radiusTrackingExitTime_le,
    AbsorptionCutoff.isStoppingTime_radiusTrackingExitTime,
    AbsorptionCutoff.radiusDifferenceQuotient_mul_trackingError,
    AbsorptionCutoff.radiusTrackingError_zero_of_eval_zero_eq,
    AbsorptionCutoff.radiusTrackingError_succ,
    AbsorptionCutoff.radiusTrackingError_succ_eq_mul_add,
    AbsorptionCutoff.normalizedRadiusError_succ,
    AbsorptionCutoff.markovPathMeasure_dirac_ae_eval_zero_eq,
    AbsorptionCutoff.markovPathMeasure_dirac_ae_eval_nonneg,
    AbsorptionCutoff.stoppedRadiusMultiplier,
    AbsorptionCutoff.measurable_stoppedRadiusMultiplier,
    AbsorptionCutoff.stronglyMeasurable_stoppedRadiusMultiplier_mul_normalizedRadiusError,
    AbsorptionCutoff.roundedBeta_nonneg, AbsorptionCutoff.roundedAlpha_nonneg,
    AbsorptionCutoff.abs_radiusNormalizedMultiplier_le_roundedAlpha_of_lt_exit_ae,
    AbsorptionCutoff.abs_stoppedRadiusMultiplier_le_roundedAlpha_ae,
    AbsorptionCutoff.sq_stoppedRadiusMultiplier_le_roundedAlpha_sq_ae,
    AbsorptionCutoff.stoppedNormalizedRadiusNoise,
    AbsorptionCutoff.stoppedNormalizedRadiusNoise_succ,
    AbsorptionCutoff.normalizedRadiusError_succ_eq_stopped_of_lt_exit,
    AbsorptionCutoff.stronglyMeasurable_stoppedNormalizedRadiusError,
    AbsorptionCutoff.integrable_sq_normalizedRadiusError,
    AbsorptionCutoff.integrable_sq_stoppedNormalizedRadiusError,
    AbsorptionCutoff.integrable_weighted_sq_stoppedNormalizedRadiusError,
    AbsorptionCutoff.integrable_sq_stoppedRadiusMultiplier_mul_normalizedRadiusError,
    AbsorptionCutoff.integrable_stoppedNormalizedRadiusNoise_succ,
    AbsorptionCutoff.condExp_stoppedNormalizedRadiusNoise_succ_eq_zero,
    AbsorptionCutoff.integrable_sq_stoppedNormalizedRadiusNoise_succ,
    AbsorptionCutoff.integrable_stoppedRadiusDrift_mul_stoppedNormalizedRadiusNoise_succ,
    AbsorptionCutoff.integrable_sq_stoppedRadiusDrift_add_stoppedNormalizedRadiusNoise_succ,
    AbsorptionCutoff.condExp_sq_stoppedNormalizedRadiusNoise_succ,
    AbsorptionCutoff.condExp_sq_stoppedNormalizedRadiusNoise_succ_le_of_raw,
    AbsorptionCutoff.exists_pos_condExp_sq_stoppedNormalizedRadiusNoise_succ_le,
    AbsorptionCutoff.exists_pos_uniform_condExp_sq_stoppedNormalizedRadiusNoise_succ_le,
    AbsorptionCutoff.stoppedNormalizedRadiusError_condExp_step,
    AbsorptionCutoff.exists_pos_stoppedNormalizedRadiusError_condExp_step,
    AbsorptionCutoff.exists_pos_uniform_stoppedNormalizedRadiusError_condExp_step,
    AbsorptionCutoff.exists_pos_stoppedNormalizedRadiusError_moment_recursion,
    AbsorptionCutoff.exists_pos_uniform_stoppedNormalizedRadiusError_moment_recursion,
    AbsorptionCutoff.exists_pos_stoppedNormalizedRadiusError_secondMoment_le,
    AbsorptionCutoff.exists_pos_uniform_stoppedNormalizedRadiusError_secondMoment_le,
    AbsorptionCutoff.exists_pos_radiusTrackingExitProbability_le,
    AbsorptionCutoff.exists_pos_uniform_radiusTrackingExitProbability_le}
  \leanok
  \uses{lem:subcritical-fixed-conditional-noise,
    lem:subcritical-normalized-lattice-amplification}
  Along the deterministic orbit
  $\bar h_{t+1}=V_{A,\varrho}(\bar h_t)$, define
  \[
    e_t=\mathscr H_t^{(N)}-\bar h_t,
    \qquad
    \xi_{t+1}=\mathscr H_{t+1}^{(N)}
      -V_{A,\varrho}(\mathscr H_t^{(N)}),
    \qquad
    \mathcal R_t=\frac{e_t}{\bar h_t+a}.
  \]
  If $\mathscr H_0^{(N)}=\bar h_0$, then $e_0=0$ and the difference quotient
  $D_{t,N}$ gives
  \[
    e_{t+1}=D_{t,N}e_t+\xi_{t+1},
    \qquad
    \mathcal R_{t+1}
      =(D_{t,N}\beta_{t,N})\mathcal R_t
        +\beta_{t,N}\frac{\xi_{t+1}}{\bar h_t+a}.
  \]
  Let $\tau_\delta\leq T+1$ be the first failure of the normalized tracking
  bound, with $T+1$ as sentinel, and set
  \[
    A_{t,N}:=\mathbf 1_{\{t<\tau_\delta\}}D_{t,N}\beta_{t,N},
    \qquad
    \eta_{t+1,N}:=
      \mathbf 1_{\{t<\tau_\delta\}}\beta_{t,N}
      \frac{\xi_{t+1}}{\bar h_{t,N}+a}.
  \]
  Then $A_{t,N}$ is adapted, $A_{t,N}^2\leq\alpha_{t,N}^2$ almost surely, and
  the stopped recursion has conditionally centered noise satisfying
  \[
    \overline{\mathcal R}_{t+1}
      =A_{t,N}\overline{\mathcal R}_t+\eta_{t+1,N},
    \qquad
    \mathbb E[\eta_{t+1,N}^2\mid\mathcal F_t]
      \leq \frac{C\beta_{t,N}^2}{N}
      \quad\text{a.s.}
  \]
  Consequently,
  \[
    \mathbb E[\overline{\mathcal R}_{T+1}^2]
      \leq \frac{C(T+1)}{N}
        \operatorname{Amp}_{T+1}(\alpha,\beta)^2,
    \qquad
    \mathbb P(\tau_\delta\leq T)
      \leq \frac{C(T+1)}{N\delta^2}
        \operatorname{Amp}_{T+1}(\alpha,\beta)^2.
  \]
\end{lemma}
\begin{proof}[Proof outline]
  \leanok
  \uses{lem:subcritical-fixed-conditional-noise,
    lem:subcritical-normalized-lattice-amplification}
  The difference-quotient identity gives the exact error recursion. Before
  $\tau_\delta$, localization to the rounded orbit window bounds its normalized
  multiplier by $\alpha_{t,N}$. All stopped quantities are measurable for the
  coordinate-prefix filtration, and the one-step noise remains conditionally
  centered after multiplication by the adapted stopping coefficient.

  The conditional variance estimate follows from
  \Cref{lem:subcritical-fixed-conditional-noise}. The abstract stopped-moment
  estimate, with backward weights built from $\alpha_{t,N}$, then telescopes to
  the bound at $T+1$. On $\{\tau_\delta\leq T\}$, the stopped normalized error
  exceeds $\delta$ at the exit index. Chebyshev's inequality gives the stated
  exit-probability estimate.
\end{proof}

\begin{proposition}[Fixed-precision radius concentration]
  \label{prop:subcritical-exact-radius-concentration}
  \lean{AbsorptionCutoff.subcritical_exact_radius_concentration}
  \leanok
  \uses{lem:subcritical-radius-tracking-recursion,
    lem:subcritical-normalized-lattice-amplification}
  Assume $0<A<A_{\mathrm{lat}}$, fix $\varrho\in(0,1)$, and let
  $x_N\in[-1,1]^N$.  If
  \[
    \frac{\bar T_N}{N}\exp\{C(\log\log N)^2\}\longrightarrow0
    \qquad\text{for every fixed }C\in\mathbb R,
  \]
  then, under the canonical rounded-radius path law started from the rounded
  squared radius of $x_N$,
  \[
    \sup_{0\leq t\leq\bar T_N}
    \frac{|\mathscr H_t^{(N)}-\bar h_{t,N}|}
      {\bar h_{t,N}+(\log N)^{-1}}
    \xrightarrow{\mathbb P}0.
  \]
  Formally, for every $\delta>0$, the probability that the displayed
  normalized error exceeds $\delta$ at some $t\leq\bar T_N$ tends to zero.
  The deterministic initial-radius estimate gives
  $\bar h_{0,N}\leq(\varrho^{-1}+1/2)^2$ uniformly in $N$.  The uniform
  amplification theorem and the uniform stopped-exit estimate then bound this
  probability by a constant multiple of
  \[
    \frac{\bar T_N+1}{N}
      \exp\{C'(\log\log N)^2\}.
  \]
  The hypothesis handles the $\bar T_N$ term, while
  $N^{-1}\exp\{C'(\log\log N)^2\}\to0$ handles the sentinel correction.
\end{proposition}
\begin{proof}
  \leanok
  \uses{lem:subcritical-radius-tracking-recursion,
    lem:subcritical-normalized-lattice-amplification}
\end{proof}

\begin{lemma}[Zero-bin criterion]
  \label{lem:subcritical-exact-zero-bin}
  \lean{AbsorptionCutoff.Hmap_eq_zero_iff, AbsorptionCutoff.Hmap_ne_zero_iff,
    AbsorptionCutoff.Hmap_eq_zero_iff_forall_observable,
    AbsorptionCutoff.Hmap_ne_zero_iff_exists_observable}
  \leanok
  \uses{def:Hmap, lem:Q1-zero, lem:subcritical-layer-formulas-exact}
  Fix $\varrho\in(0,1)$ and write
  $b_\varrho:=\varrho^{-1}\operatorname{arctanh}(\varrho/2)=
  b_{\varrho,0}$ (the first admissible layer threshold
  $\texttt{roundedLayerThreshold}\ \varrho\ 0$).  One rounded Gaussian step lands in
  the absorbing zero bin, $H_{A,\varrho,N}(h,g)=0$, exactly when every rescaled
  coordinate stays within that threshold:
  \[
    H_{A,\varrho,N}(h,g)=0
    \iff
    \max_{1\leq i\leq N}|A\sqrt h\,g_i|\leq b_\varrho .
  \]
  Indeed, the nonnegative empirical average vanishes iff each squared rounded
  coordinate does, i.e.\ $Q_1(\varrho^{-1}\tanh(\varrho A\sqrt h\,g_i))=0$; by
  \Cref{lem:Q1-zero} this is
  $|\varrho^{-1}\tanh(\varrho A\sqrt h\,g_i)|\leq\tfrac12$, and the monotone layer
  preimage of \Cref{lem:subcritical-layer-formulas-exact} converts it to
  $|A\sqrt h\,g_i|\leq b_\varrho$.  Negating gives the survival form
  ($\texttt{Hmap\_ne\_zero\_iff}$): $H_{A,\varrho,N}(h,g)\neq0$ iff some coordinate
  satisfies $b_\varrho<|A\sqrt h\,g_i|$.
\end{lemma}
\begin{proof}\leanok\uses{lem:Q1-zero, lem:subcritical-layer-formulas-exact}\end{proof}

\begin{lemma}[One-step escape probability]
  \label{lem:subcritical-exact-escape-probability}
  \lean{AbsorptionCutoff.measureReal_Hmap_ne_zero_le,
    AbsorptionCutoff.measureReal_Hkernel_compl_singleton_zero_le,
    AbsorptionCutoff.tendsto_escape_at_eta_atTop}
  \leanok
  \uses{lem:subcritical-exact-zero-bin, lem:subcritical-fixed-moment-bound, def:Hkernel}
  Fix $0<A$, $\varrho\in(0,1)$, and $h>0$.  Under the product Gaussian
  $\mathcal G_N$, the probability that one rounded step escapes the absorbing zero
  bin obeys the Gaussian union bound
  \[
    \mathcal G_N\{g : H_{A,\varrho,N}(h,g)\neq0\}
    \leq
    2N\exp\!\left(
      -\tfrac12\Big(\tfrac{b_\varrho}{A\sqrt h}\Big)^2
    \right).
  \]
  By the survival criterion the escape event is
  $\bigcup_{i<N}\{g_i:U_{A,\varrho}(h,g_i)\neq0\}$; each coordinate marginal of
  $\mathcal G_N$ is the standard Gaussian, so the union bound and the one-coordinate
  support estimate of \Cref{lem:subcritical-fixed-moment-bound} give the claim.
  Transferring through the Gaussian pushforward $H_{A,\varrho,N}(h,\cdot)_\#\mathcal
  G_N$ (\Cref{def:Hkernel}) gives the kernel form: the non-absorption mass
  $H_{A,\varrho,N}(h,\{0\}^c)$ satisfies the same bound.  Finally, evaluating at the
  terminal radius scale $h=\eta/\log N$ with $0<\eta<b_\varrho^2/(2A^2)$, the bound
  equals $2N^{1-b_\varrho^2/(2A^2\eta)}$ and tends to $0$ (the exponent is negative) --
  the paper's final-absorption rate.
\end{lemma}
\begin{proof}\leanok
  \uses{lem:subcritical-exact-zero-bin, lem:subcritical-fixed-moment-bound}
\end{proof}

\begin{lemma}[Terminal small-radius decay]
  \label{lem:subcritical-exact-terminal-decay}
  \lean{AbsorptionCutoff.exists_roundedMeanMap_le_rpow_atTop,
    AbsorptionCutoff.tendsto_roundedMeanMap_mul_log_of_le_fixedPrecisionScale,
    AbsorptionCutoff.fixedPrecisionScale_mul_log,
    AbsorptionCutoff.tendsto_add_fixedPrecisionScale_mul_log_of_le,
    AbsorptionCutoff.tendsto_roundedOrbit_terminal_mul_log,
    AbsorptionCutoff.tendsto_roundedOrbit_terminal_add_scale_mul_log}
  \leanok
  \uses{lem:subcritical-fixed-precision-map-estimates,
    lem:subcritical-normalized-lattice-amplification}
  There are constants $c,C>0$ such that, for all sufficiently large $N$, every radius
  $h$ at or below the terminal scale $a_N=1/\log N$ already satisfies
  \[
    V_{A,\varrho}(h)\leq C\,N^{-c}.
  \]
  Compose the small-radius map bound $V_{A,\varrho}(h)\leq C\exp(-c/h)$
  (\Cref{lem:subcritical-fixed-precision-map-estimates}) with $1/h\geq\log N$, using
  $\exp(-c\log N)=N^{-c}$; the case $h=0$ is immediate since $V_{A,\varrho}(0)=0$.  This
  is the deterministic input behind the paper's terminal-entrance estimate
  $\bar h_{s_N,N}\leq C N^{-c}$.  Consequently, for any pre-terminal radius sequence
  $h_N$ with $0\leq h_N\leq a_N$ eventually, $V_{A,\varrho}(h_N)\log N\to0$ (since
  $N^{-c}\log N\to0$); this is the paper's $\bar h_{s_N,N}\log N\to0$.  Since moreover
  $a_N\log N=1$ exactly, the normalized denominator satisfies
  $(V_{A,\varrho}(h_N)+a_N)\log N\to1$, the paper's $(\bar h_{s_N,N}+a_N)\log N\to1$.
  Specializing $h_N$ to the deterministic orbit at its $a_N$-entrance
  ($\bar t_N=\texttt{roundedOrbitEntrance}$, using the entrance bound
  $\bar h_{\bar t_N,N}\leq a_N$) yields the paper's $\bar h_{s_N,N}\log N\to0$ with
  $s_N=\bar t_N+1$.
\end{lemma}
\begin{proof}\leanok\uses{lem:subcritical-fixed-precision-map-estimates}\end{proof}

\begin{theorem}[Fixed-precision subcritical dimension cutoff]
  \label{thm:subcritical-dimension-cutoff}
  \lean{AbsorptionCutoff.subcritical_dimension_cutoff,
    AbsorptionCutoff.subcritical_dimension_cutoff_roundedVector_paper,
    AbsorptionCutoff.subcritical_dimension_cutoff_mixingTime_bounds,
    AbsorptionCutoff.subcritical_dimension_cutoff_roundedVector,
    AbsorptionCutoff.subcritical_dimension_cutoff_roundedVector_mixingTime_bounds}
  \leanok
  \uses{prop:subcritical-exact-radius-concentration,
    lem:subcritical-exact-zero-bin,
    lem:subcritical-exact-escape-probability,
    lem:subcritical-exact-terminal-decay,
    lem:Hkernel-tv}
  Assume $0<A<A_{\mathrm{lat}}$, fix $\varrho\in(0,1)$, and let
  $x_N\in[-1,1]^N$ be deterministic.  Write $\bar t_N$ for the first time at
  which the rounded deterministic radius orbit reaches
  $a_N=(\log N)^{-1}$, and suppose $\bar t_N\to\infty$.  Then, under the
  canonical rounded-radius path law,
  \[
    \mathbb P(\tau_\varrho^{(N)}>\bar t_N-1)\longrightarrow1,
    \qquad
    \mathbb P(\tau_\varrho^{(N)}>\bar t_N+2)\longrightarrow0.
  \]
  Equivalently, the total-variation distances of both the scalar
  squared-radius laws and the rounded vector laws
  $P_{\varrho,A,N}^t(Q_\varrho x_N,\cdot)$ from their respective point masses
  at zero tend respectively to $1$ and $0$.
  Consequently, for every fixed $\varepsilon\in(0,1)$, eventually
  \[
    \bar t_N-1<t_{\mathrm{mix}}^{(N)}(\varepsilon)
      \leq \bar t_N+2,
  \]
  and hence $t_{\mathrm{mix}}^{(N)}(\varepsilon)=\bar t_N+O(1)$, for both the
  scalar and rounded vector chains.
\end{theorem}
\begin{proof}
  The geometric deterministic contraction gives
  $\bar t_N+1=O(\log\log N)$, hence the exact rate hypothesis of
  \Cref{prop:subcritical-exact-radius-concentration}.  One step past the
  terminal entrance, the deterministic radius times $\log N$ tends to zero
  and its normalizing denominator tends to one, so concentration yields
  $\mathscr H_{\bar t_N+1}^{(N)}\log N\to0$ in probability.
  Splitting one further kernel step at $\eta/\log N$ and applying
  \Cref{lem:subcritical-exact-escape-probability} proves the upper limit.
  For the lower limit, the deterministic orbit at $\bar t_N-1$ is strictly
  above $a_N$; an exact radius equal to zero therefore forces normalized
  tracking error greater than $1/2$, whose probability vanishes by
  concentration.  The path-space survival and total-variation forms follow
  from \Cref{lem:Hkernel-tv}.
  \leanok
\end{proof}

\begin{lemma}[Compact inward persistence margin]
  \label{lem:metastable-inward-margin}
  \lean{AbsorptionCutoff.exists_roundedPositiveDriftComponent_inward_margin,
    AbsorptionCutoff.metastableExitEvent,
    AbsorptionCutoff.measurableSet_metastableExitEvent,
    AbsorptionCutoff.one_sub_measureReal_metastableExitEvent_le_survival,
    AbsorptionCutoff.exists_preexit_step_deviation_gt_half,
    AbsorptionCutoff.shiftedFiniteHorizonStepDeviationEvent,
    AbsorptionCutoff.markovPathMeasure_measureReal_shiftedFiniteHorizonStepDeviationEvent_le,
    AbsorptionCutoff.markovPathMeasure_measureReal_inter_start_metastableExitEvent_le,
    AbsorptionCutoff.exists_uniform_markovPathMeasure_exit_exp_bound,
    AbsorptionCutoff.tendsto_one_add_floor_exp_mul_exp_neg,
    AbsorptionCutoff.exists_exponential_absorption_survival}
  \leanok
  Let $(h_-,h_+)$ be the rightmost positive-drift component of the rounded
  mean map.  There is $\eta_0>0$ such that, whenever $0<\eta<\eta_0$ and
  $J_\eta=[h_+-\eta,h_++\eta]$,
  \[
    V_{A,\varrho}(J_\eta)\subset J_\eta^\circ.
  \]
  More precisely, there is $d_\eta>0$ such that
  \[
    V_{A,\varrho}(h)\in
    [h_+-\eta+d_\eta,h_++\eta-d_\eta]
    \qquad (h\in J_\eta).
  \]
  There are also $T_\eta,C,c_0>0$ such that
  \[
    \mathbb P\!\left(
      \max_{0\leq s\leq T}
      |\mathscr H_{T_\eta+s}^{(N)}-h_+|>\eta
    \right)
    \leq C(1+T)e^{-c_0N}.
  \]
  Therefore, for $0<c_1<c_0$ and $T_N=\lfloor e^{c_1N}\rfloor$,
  \[
    \mathbb P(\tau>T_\eta+T_N)\longrightarrow1,
  \]
  uniformly for initial states in the prescribed compact set.
\end{lemma}
\begin{proof}[Proof outline]
  \leanok
  Positive drift at the left endpoint, right stability at the right endpoint,
  and monotonicity give strict inward motion on $J_\eta$. Compactness supplies
  the margin $d_\eta$.

  If a path enters the half-sized well and later exits the full well, some
  genuinely pre-exit step deviates from the mean map by more than
  $d_\eta/2$. The one-step Hoeffding estimate and a stopped union bound give
  an error of order $Te^{-cN}$. Combine this with finite-horizon stochastic
  entrance. For $T=T_N$, the error tends to zero because $c_1<c_0$; remaining
  in the well forces survival beyond $T_\eta+T_N$.
\end{proof}

\begin{lemma}[Uniform one-step absorption chance]
  \label{lem:uniform-rounded-absorption-chance}
  \lean{AbsorptionCutoff.exists_pos_le_measureReal_Hkernel_singleton_zero,
    AbsorptionCutoff.exists_geometric_Hkernel_survival_bound,
    AbsorptionCutoff.tendsto_Hkernel_survival_and_ae_absorption}
  \leanok
  For fixed $N$, there is $p_N>0$ such that every starting radius in the
  canonical interval $[0,M_\varrho]$ hits zero in one step with probability
  at least $p_N$. Indeed, a fixed coordinatewise Gaussian box of positive
  product measure makes every rounded coordinate vanish, uniformly over the
  starting radius. Absorption at zero then yields the geometric survival bound
  \[
    H_{A,\varrho,N}^t(q,\{0\}^{\mathsf c})\leq(1-p_N)^t.
  \]
  Hence survival tends to zero, and the event that the canonical absorption
  time equals infinity has measure zero.
\end{lemma}
\begin{proof}\leanok\end{proof}

\begin{theorem}[Rounded qualitative metastability and absorption]
  \label{thm:rounded-qualitative-metastability}
  \lean{AbsorptionCutoff.rounded_qualitative_metastability_paper,
    AbsorptionCutoff.ae_absorption_roundedPkernel_of_rounded_state}
  \leanok
  \uses{lem:metastable-inward-margin,lem:uniform-rounded-absorption-chance}
  For positive dimensions, the rounded chain remains in a compact subset of the
  rightmost positive-drift component for an exponentially long time with
  probability tending to one as the dimension grows, uniformly over the
  prescribed compact initial set.
  In each fixed dimension, from every initial state in the full finite rounded
  state space
  \[
    (\varrho\mathbb Z)^N\cap
    [-1-\varrho/2,1+\varrho/2]^N,
  \]
  the same chain is absorbed at zero almost surely, and its survival probability
  tends to zero. These are the two clauses of the paper's qualitative
  metastability theorem.
\end{theorem}
\begin{proof}\leanok\end{proof}

\begin{lemma}[Deterministic affine-recursion entrance scaffold]
  \label{lem:affine-entrance-scaffold}
  \lean{AbsorptionCutoff.affineRecursion, AbsorptionCutoff.affineEntranceTime,
    AbsorptionCutoff.lt_affineEntranceTime_iff,
    AbsorptionCutoff.affineSurvivalSet_eq_lt_affineEntranceTime,
    AbsorptionCutoff.affineRecursion_succ_le_mul_add, AbsorptionCutoff.affineRecursion_le_mul_prod,
    AbsorptionCutoff.neg_log_div_lt_sum_log_of_affineRecursion,
    AbsorptionCutoff.measure_sum_ge_le_exp_mul_prod_mgf,
    AbsorptionCutoff.measure_sum_ge_le_exp_mul_mgf_pow, AbsorptionCutoff.affineSurvivalSet,
    AbsorptionCutoff.affineSurvivalSet_subset_log_sum, AbsorptionCutoff.logAffineIncrement,
    AbsorptionCutoff.measurable_logAffineIncrement, AbsorptionCutoff.iIndepFun_logAffineIncrement,
    AbsorptionCutoff.identDistrib_logAffineIncrement,
    AbsorptionCutoff.integrable_exp_mul_sum_logAffineIncrement,
    AbsorptionCutoff.measurable_affineRecursion,
    AbsorptionCutoff.measurableSet_affineSurvivalSet,
    AbsorptionCutoff.measureReal_affineSurvivalSet_le,
    AbsorptionCutoff.measureReal_affineSurvivalSet_le_exp_sub,
    AbsorptionCutoff.measureReal_affineSurvivalSet_le_exp_mul_exp_neg,
    AbsorptionCutoff.affineThreshold_sub_one_lt_floor,
    AbsorptionCutoff.measureReal_affineSurvivalSet_floor_le,
    AbsorptionCutoff.exists_pos_integral_log_add_neg_of_tendsto,
    AbsorptionCutoff.tendsto_integral_log_add_inv_nat,
    AbsorptionCutoff.exists_pos_integral_log_add_neg,
    AbsorptionCutoff.exists_pos_lt_one_of_hasDerivAt_neg,
    AbsorptionCutoff.hasDerivAt_mgf_zero_integral,
    AbsorptionCutoff.exists_pos_mgf_lt_one,
    AbsorptionCutoff.exists_pos_exp_neg_ge,
    AbsorptionCutoff.exists_pos_mgf_le_exp_neg,
    AbsorptionCutoff.exists_pos_mem_lt_one_of_hasDerivAt_neg,
    AbsorptionCutoff.exists_pos_integrable_mgf_lt_one,
    AbsorptionCutoff.exists_pos_integrable_mgf_le_exp_neg,
    AbsorptionCutoff.measureReal_affineEntranceTime_gt_floor_le}
  \leanok
  Let
  \[
    U_0=K,
    \qquad
    U_{t+1}=M_{t+1}U_t+b,
    \qquad
    T=\inf\{t:U_t\leq K_\ast\}.
  \]
  If $b\leq\varepsilon K_\ast$, then on $\{T>n\}$,
  \[
    U_n\leq K\prod_{j<n}(M_{j+1}+\varepsilon),
    \qquad
    \sum_{j<n}\log(M_{j+1}+\varepsilon)
      >-\log(K/K_\ast).
  \]
  Suppose that the shifted logarithmic increments are independent and
  identically distributed, have negative mean, and admit an exponential
  moment near the origin. Then there exist $s,c_3>0$ such that
  $e^{sX_1}$ is integrable and
  $\operatorname{mgf}_{X_1}(s)\leq e^{-c_3}$. With
  \[
    c_1=s/c_3,
    \qquad
    c_2=\max\{e^{-s\log K_\ast+c_3},1\},
  \]
  one has, uniformly in $K$ and $r\geq0$,
  \[
    \mathbb P\!\left(T>\lfloor c_1\log K+r\rfloor\right)
      \leq c_2e^{-c_3r}.
  \]
\end{lemma}
\begin{proof}[Proof outline]
  \leanok
  Iterating $U_{t+1}\leq(M_{t+1}+\varepsilon)U_t$ before $T$ gives the
  pathwise product bound and the logarithmic event inclusion. Independence
  factors the exponential moment of the finite sum, and the iid Chernoff
  bound reduces it to the $n$-th power of one moment-generating function.

  Dominated convergence selects $\varepsilon>0$ with
  $\Ee\log(M_1+\varepsilon)<0$. The derivative of the moment-generating
  function at zero is this negative mean, so an integrable parameter $s>0$
  may be chosen with subunit mgf. Writing this value as at most $e^{-c_3}$
  and choosing the displayed $c_1,c_2$ gives the floored survival estimate.
\end{proof}

\begin{lemma}[Logarithmic entrance for contractive affine recursions]
  \label{lem:negative-drift-affine-entrance}
  \lean{AbsorptionCutoff.exists_affineEntrance_bound}
  \leanok
  \uses{lem:affine-entrance-scaffold}
  Let $(M_j)_{j\ge1}$ be independent, identically distributed, nonnegative multipliers with
  $\Ee\log M_1<0$, and assume that for every $\varepsilon>0$ the shifted increment
  $\log(M_1+\varepsilon)$ has an exponential moment in a neighborhood of the origin. For every
  $b$ there exist $K_\ast,c_1,c_2,c_3>0$, depending only on $b$ and the law of $M_1$ (uniform in
  the start $K$ and excess $r$), such that the affine recursion $U_0=K$,
  $U_{t+1}=M_{t+1}U_t+b$ satisfies
  \[
    \mathbb P\!\left(\inf\{t:U_t\le K_\ast\} > \lfloor c_1\log K + r\rfloor\right)
    \le c_2\,e^{-c_3 r}.
  \]
  This is \texttt{eq:negative-drift-affine-entrance}: choose $\varepsilon$ with
  $\Ee\log(M_1+\varepsilon)<0$, set $K_\ast=\max(b/\varepsilon,1)$ (so $b\le\varepsilon K_\ast$),
  and apply the scaffold's parametrized capstone uniformly in $K,r$. Since the entrance time
  is $\mathbb N\cup\{\infty\}$-valued, the floored threshold $\lfloor c_1\log K+r\rfloor$ agrees
  with the paper's real threshold $c_1\log K+r$.
\end{lemma}
\begin{proof}\leanok\uses{lem:affine-entrance-scaffold}\end{proof}

\section{Reduction to the squared-radius chain}
\label{sec:vector-radius-reduction}

This section begins the paper's \texttt{prop:gaussian-tv-reduction}: after one Gaussian
step the vector law depends on the previous state only through its normalized squared
radius $r_N(x)=N^{-1}\|x\|_2^2$, factorizing the vector kernel $P_{A,N}$ through the scalar
squared-radius kernel $K_{A,N}$ and a reconstruction kernel $J_{A,N}$.

\begin{lemma}[Gaussian radial Gamma law and negative moments]
  \label{lem:gaussian-squared-norm-gamma}
  \lean{AbsorptionCutoff.gaussianSquaredNorm, AbsorptionCutoff.gaussianEuclideanNorm,
    AbsorptionCutoff.map_gaussianSquaredNorm_gaussianVec,
    AbsorptionCutoff.mgf_id_gammaMeasure_of_lt,
    AbsorptionCutoff.map_eq_of_mgf_eq_on_Ioo,
    AbsorptionCutoff.map_gaussianSquaredNorm_eq_gammaMeasure,
    AbsorptionCutoff.integral_neg_rpow_gaussianSquaredNorm,
    AbsorptionCutoff.integral_neg_rpow_gaussianEuclideanNorm,
    AbsorptionCutoff.integral_neg_two_gaussianEuclideanNorm,
    AbsorptionCutoff.integral_neg_two_normalized_gaussianEuclideanNorm,
    AbsorptionCutoff.integral_neg_two_normalized_gaussianEuclideanNorm_le_two,
    AbsorptionCutoff.integral_neg_two_scaled_gaussianEuclideanNorm,
    AbsorptionCutoff.integral_neg_two_scaled_gaussianEuclideanNorm_le,
    AbsorptionCutoff.integral_neg_two_scaled_gaussianEuclideanNorm_lt_one_of_dimension,
    AbsorptionCutoff.exists_dimension_threshold_for_reciprocal_moment,
    AbsorptionCutoff.eventually_integral_neg_two_scaled_gaussianEuclideanNorm_lt_one,
    AbsorptionCutoff.Fmap_div_tendsto_gaussianSquaredNorm,
    AbsorptionCutoff.antitoneOn_Fmap_div,
    AbsorptionCutoff.Fmap_div_ge_of_pos_of_le,
    AbsorptionCutoff.tanh_one_sq_mul_min_sq_one_le,
    AbsorptionCutoff.sum_tanh_sq_lower_gaussianSquaredNorm,
    AbsorptionCutoff.inv_Fmap_div_le_one_add_inv_gaussianSquaredNorm,
    AbsorptionCutoff.integrable_inv_Fmap_div,
    AbsorptionCutoff.integrable_inv_Kchain,
    AbsorptionCutoff.exists_integral_inv_Fmap_le_on_Icc,
    AbsorptionCutoff.integral_inv_Fmap_div_tendsto,
    AbsorptionCutoff.exists_small_radius_integral_inv_Fmap_div_lt,
    AbsorptionCutoff.exists_small_radius_integral_inv_Fmap_div_lt_one,
    AbsorptionCutoff.exists_Kchain_inv_foster_bound,
    AbsorptionCutoff.exists_uniform_integral_inv_Kchain_pow_le,
    AbsorptionCutoff.exists_uniform_integral_inv_cesaroMeasure_le}
  \leanok
  For $N>0$, the squared norm of a standard Gaussian vector in $\mathbb R^N$
  has Gamma law
  \[
    (\texttt{gaussianSquaredNorm}\ N)_\#\mathcal G_N
    =\operatorname{Gamma}(N/2,1/2).
  \]
  Hence, for $p<N/2$ and $q<N$,
  \[
    \mathbb E[(\chi_N^2)^{-p}]
      =(1/2)^p\frac{\Gamma(N/2-p)}{\Gamma(N/2)},
    \qquad
    \mathbb E[\|G\|_2^{-q}]
      =(1/2)^{q/2}\frac{\Gamma((N-q)/2)}{\Gamma(N/2)}.
  \]
  In particular, if $N>2$ and $A\ne0$, then
  \[
    \mathbb E[\|G\|_2^{-2}]=\frac1{N-2},
    \qquad
    \lim_{q\downarrow0}
    \mathbb E\!\left[\left(\frac{F_{A,N}(q,G)}q\right)^{-1}\right]
    =\frac{N}{A^2(N-2)}.
  \]
  If $2A^2<(A^2-1)N$, there are $a<1$ and $B<\infty$ such that
  \[
    K_{A,N}(q,y^{-1})\le a q^{-1}+B,\qquad 0<q\le1.
  \]
  Consequently the reciprocal radius is integrable, with a bound uniform over
  all iterated laws and positive-length Cesàro averages started from
  $q\in(0,1]$.
\end{lemma}
\begin{proof}[Proof outline]
  \leanok
  The one-coordinate squared-Gaussian and Gamma moment-generating functions,
  independence, and uniqueness of the complex MGF identify the Gamma law. The
  density integral gives the two negative-moment formulas.

  The map $q\mapsto F_{A,N}(q,g)/q$ is nonincreasing, and
  \[
    \tanh^2(1)\min\{A^2R,1\}\frac{S}{1+S}
      \leq \sum_i\tanh^2(A\sqrt R\,g_i),
    \qquad S=\|g\|_2^2,
  \]
  provides an integrable reciprocal dominator when $N>2$. Dominated
  convergence gives the limit at zero. Under the dimension criterion this
  limit is below one, giving the drift estimate near zero; compactness gives
  the finite bound away from zero. Iterating the affine drift inequality
  yields the uniform power and Cesàro bounds.
\end{proof}

\begin{lemma}[One-coordinate truncated Gaussian Laplace decay]
  \label{lem:gaussian-truncated-coordinate-laplace}
  \lean{AbsorptionCutoff.truncatedGaussianSquare,
    AbsorptionCutoff.measurable_truncatedGaussianSquare,
    AbsorptionCutoff.truncatedGaussianSquare_nonneg,
    AbsorptionCutoff.tendsto_integral_truncatedGaussianSquare_atTop,
    AbsorptionCutoff.exists_one_le_integral_truncatedGaussianSquare_gt_inv_sq,
    AbsorptionCutoff.exists_truncatedGaussianSquare_parameters,
    AbsorptionCutoff.truncatedGaussianSquareMean,
    AbsorptionCutoff.integral_gaussianVec_truncatedGaussianSquare_eval_eq_mean,
    AbsorptionCutoff.hasSubgaussianMGF_truncatedGaussianSquare_eval_sub_mean,
    AbsorptionCutoff.iIndepFun_truncatedGaussianSquare_eval_sub_mean,
    AbsorptionCutoff.measureReal_truncatedGaussianAverage_le_mean_sub_le,
    AbsorptionCutoff.truncatedGaussianSum,
    AbsorptionCutoff.truncatedGaussianAverage,
    AbsorptionCutoff.measurable_truncatedGaussianSum,
    AbsorptionCutoff.truncatedGaussianSum_nonneg,
    AbsorptionCutoff.measurable_truncatedGaussianAverage,
    AbsorptionCutoff.truncatedGaussianAverage_nonneg,
    AbsorptionCutoff.ae_truncatedGaussianAverage_one_pos,
    AbsorptionCutoff.iIndepFun_gaussian_coordinate_truncatedSquare,
    AbsorptionCutoff.integral_exp_neg_truncatedGaussianSum,
    AbsorptionCutoff.integral_exp_neg_truncatedGaussianSum_one_le,
    AbsorptionCutoff.integral_exp_neg_truncatedGaussianSquare_one_le_add,
    AbsorptionCutoff.integral_exp_neg_truncatedGaussianSquare_one_le}
  \leanok
  \uses{lem:gaussian-squared-norm-gamma}
  For a standard Gaussian coordinate $G$, set
  $U_L=\min\{G^2,L^2\}$. Then $\mathbb E U_L\to1$ as $L\to\infty$, and for
  $t\geq1$,
  \[
    \mathbb E e^{-tU_1}\leq \frac2{\sqrt t}.
  \]
  For every $A>1$, one may therefore choose $L\geq1$,
  $\varepsilon\in(0,1)$ and $\delta>0$ so that
  $(1-\varepsilon)A^2(\mathbb E U_L-\delta)>1$.
  If $U_{L,1},\ldots,U_{L,N}$ are independent copies and
  $\mu_L=\mathbb E U_L$, then
  \[
    \mathbb P\!\left(
      \frac1N\sum_{i=1}^N U_{L,i}\le\mu_L-\varepsilon
    \right)
    \le \exp\!\left(-\frac{2N\varepsilon^2}{L^4}\right).
  \]
  Moreover, for $t\geq1$,
  \[
    \mathbb E\exp\!\left(-t\sum_{i=1}^N\min\{G_i^2,1\}\right)
    \le \left(\frac{2}{\sqrt t}\right)^N.
  \]
\end{lemma}
\begin{proof}[Proof outline]
  \leanok
  \uses{lem:gaussian-squared-norm-gamma}
  Split the one-coordinate integral at $G^2=1$ and use the exact Gaussian
  square transform to obtain the Laplace bound. Dominated convergence gives
  $\mathbb E U_L\to1$. Hoeffding's inequality applies because
  $0\leq U_L\leq L^2$, while independence factors the finite-dimensional
  Laplace transform into the $N$-th power of its one-coordinate counterpart.
\end{proof}

\begin{lemma}[Truncated Gaussian negative moments]
  \label{lem:gaussian-truncated-negative-moment}
  \lean{AbsorptionCutoff.integrable_neg_rpow_truncatedGaussianAverage_one_and_integral_le,
    AbsorptionCutoff.integrable_neg_rpow_truncatedGaussianAverage_one_and_integral_le_exp,
    AbsorptionCutoff.integrable_neg_rpow_truncatedGaussianAverage_one_mul_dimension_and_integral_le,
    AbsorptionCutoff.integrable_neg_rpow_truncatedGaussianAverage_mul_dimension_and_integral_le}
  \leanok
  \uses{lem:gaussian-truncated-coordinate-laplace,
    lem:laplace-envelope-negative-moment}
  Let $S_N=N^{-1}\sum_{i=1}^N\min\{G_i^2,1\}$. For
  $0<p<N/2$, the random variable $S_N^{-p}$ is integrable and
  \[
    \mathbb E S_N^{-p}
    \leq \Gamma(p)^{-1}
      \left[
        \frac{N^p}{p}
        +2^N N^{N/2}\frac{N^{p-N/2}}{N/2-p}
      \right].
  \]
  Equivalently,
  \[
    \mathbb E S_N^{-p}
    \le e^{p+1}(N/p)^p
      \left(1+\frac{p\,2^N}{N/2-p}\right).
  \]
  In particular, for $0<\alpha<1/2$,
  \[
    \mathbb E S_N^{-\alpha N}
    \le \frac{e}{1-2\alpha}
      \left(2e^\alpha\alpha^{-\alpha}\right)^N.
  \]
  These bounds remain valid when the truncation level $1$ is replaced by any
  $L\geq1$.
\end{lemma}
\begin{proof}[Proof outline]
  \leanok
  \uses{lem:gaussian-truncated-coordinate-laplace,
    lem:laplace-envelope-negative-moment}
  Split the Mellin integral at $t=N$. Below the breakpoint the Laplace
  transform is at most one; above it, tensorization gives the integrable
  power-law envelope. The Gamma quotient bound yields the exponential form.
  Increasing $L$ only increases the empirical average, so its negative power
  decreases.
\end{proof}

\begin{lemma}[Truncated Gaussian bad-event negative moment]
  \label{lem:gaussian-truncated-bad-event-moment}
  \lean{AbsorptionCutoff.integrableOn_neg_rpow_truncatedGaussianAverage_badEvent_and_integral_le,
    AbsorptionCutoff.integrableOn_neg_rpow_scaled_truncatedGaussianAverage_badEvent_and_integral_le,
    AbsorptionCutoff.integrableOn_neg_rpow_scaled_truncatedGaussianAverage_goodEvent_and_integral_le,
    AbsorptionCutoff.integrable_neg_rpow_scaled_truncatedGaussianAverage_and_integral_le,
    AbsorptionCutoff.exists_fractional_exponent_and_rate_for_scaled_truncatedGaussianAverage,
    AbsorptionCutoff.exists_eventually_integrable_neg_rpow_scaled_truncatedGaussianAverage,
    AbsorptionCutoff.exists_truncatedGaussianAverage_cramer_bound}
  \leanok
  \uses{lem:gaussian-truncated-coordinate-laplace,
    lem:gaussian-truncated-negative-moment,
    lem:set-negative-moment-holder}
  Let $0<\gamma<\alpha<1/2$, $L\ge1$, $\delta\ge0$, and $N>0$.
  On the bad event
  $S=\{\overline W_{L,N}\le\mu_L-\delta\}$, the lower negative moment is
  integrable and satisfies
  \[
    \mathbb E\!\left[
      \overline W_{L,N}^{-\gamma N}\mathbf 1_S
    \right]
    \le
    \left(\frac{e}{1-2\alpha}\right)^{\gamma/\alpha}
    \left[
      \left(2e^\alpha\alpha^{-\alpha}\right)^{\gamma/\alpha}
      \exp\!\left(
        -\frac{2\delta^2}{L^4}
        \left(1-\frac{\gamma}{\alpha}\right)
      \right)
    \right]^N.
  \]
  Consequently, for every $A>1$ there are
  $\gamma\in(0,1/2)$, $\kappa>0$, $L\ge1$, and
  $\varepsilon\in(0,1)$ such that, for all sufficiently large $N$,
  \[
    \mathbb E\!\left[
      \bigl((1-\varepsilon)A^2\overline W_{L,N}\bigr)^{-\gamma N}
    \right]
    \leq e^{-\kappa N}.
  \]
\end{lemma}
\begin{proof}[Proof outline]
  \leanok
  \uses{lem:gaussian-truncated-coordinate-laplace,
    lem:gaussian-truncated-negative-moment,
    lem:set-negative-moment-holder}
  Hölder interpolation combines the higher negative moment with the
  exponential lower-tail estimate to give the displayed bad-event bound. On
  the complementary event, the scaled average is bounded below by
  $(1-\varepsilon)A^2(\mu_L-\delta)>1$.

  The bad-event exponential base is continuous in $\gamma$ and is below one
  at $\gamma=0$; the good-event base is below one for every $\gamma>0$.
  Choose $\gamma$ small enough that both decay exponentially, then absorb the
  dimension-free prefactor by reducing the rate. Taking $\alpha=1/4$ and the
  truncation parameters from $A>1$ gives the final estimate.
\end{proof}

\begin{lemma}[Uniform small-argument hyperbolic tangent linearization]
  \label{lem:gaussian-uniform-small-tanh-ratio}
  \lean{AbsorptionCutoff.exists_pos_forall_tanh_div_self_sq_ge_one_sub,
    AbsorptionCutoff.exists_radius_truncatedGaussianSquare_le_tanh_sq_div,
    AbsorptionCutoff.scaled_truncatedGaussianAverage_le_Fmap_div_of_forall,
    AbsorptionCutoff.exists_radius_scaled_truncatedGaussianAverage_le_Fmap_div,
    AbsorptionCutoff.exists_radius_eventually_integrable_neg_rpow_Fmap_div,
    AbsorptionCutoff.exists_radius_lt_eventually_integrable_neg_rpow_Fmap_div,
    AbsorptionCutoff.exists_radius_below_stableInterval_eventually_integrable_neg_rpow_Fmap_div}
  \leanok
  \uses{lem:gaussian-truncated-bad-event-moment}
  For every $\varepsilon\in(0,1)$ there is $\eta>0$ such that
  \[
    1-\varepsilon
    \leq \left(\frac{\tanh x}{x}\right)^2
  \]
  whenever $0<|x|\leq\eta$. For fixed $A>1$ and $L\ge1$, set
  $R_0=\min\{1,\eta/(AL)\}^2$. Then, for every $g\in\mathbb R$,
  \[
    (1-\varepsilon)A^2\min\{g^2,L^2\}
    \leq \frac{\tanh^2(A\sqrt{R_0}\,g)}{R_0}
  \]
  and hence
  \[
    (1-\varepsilon)A^2\overline W_{L,N}(g)
    \leq \frac{F_{A,N}(R_0,g)}{R_0}.
  \]
  Consequently there are constants
  $\gamma\in(0,1/2)$ and $\kappa>0$ such that, for all sufficiently large
  $N$ and uniformly over $0<q\leq R_0$,
  \[
    \mathbb E\left[
      \left(\frac{F_{A,N}(q,G)}q\right)^{-\gamma N}
    \right]\leq e^{-\kappa N}.
  \]
\end{lemma}
\begin{proof}[Proof outline]
  \leanok
  The first estimate follows uniformly from $\tanh x/x\to1$ near zero. Inside
  $|g|\leq L$, it gives the coordinatewise bound;
  outside this interval, monotonicity of $\tanh$ in $|g|$ gives the same
  conclusion. Sum over coordinates and apply the truncated Cramér estimate.
  Finally use that
  $q\mapsto F_{A,N}(q,g)/q$ is nonincreasing.
\end{proof}

\begin{lemma}[Negative-moment bound away from zero]
  \label{lem:gaussian-outside-coordinate-lower}
  \lean{AbsorptionCutoff.exists_pos_forall_truncatedGaussianSquare_one_le_tanh_sq,
    AbsorptionCutoff.exists_pos_forall_scaled_truncatedGaussianAverage_one_le_Fmap,
    AbsorptionCutoff.exists_pos_forall_integrable_neg_rpow_Fmap_and_integral_le,
    AbsorptionCutoff.exists_forall_integrable_neg_rpow_Fmap_and_integral_le_exp,
    AbsorptionCutoff.integrable_rpow_Fmap_and_integral_eq_mul}
  \leanok
  \uses{lem:gaussian-squared-norm-gamma}
  For every $A>1$ and $R_0\in(0,1)$ there is $c_{A,R_0}>0$ such that,
  uniformly over $q\in[R_0,1]$ and $g\in\mathbb R$,
  \[
    c_{A,R_0}\min\{g^2,1\}
    \leq \tanh^2(A\sqrt q\,g).
  \]
  One may take
  $c_{A,R_0}=\tanh^2(1)\min\{A^2R_0,1\}$; the global two-regime
  hyperbolic-tangent bound and $q\geq R_0$ give the claim.
  Summing and normalizing gives
  $c_{A,R_0}\overline W_{1,N}\leq F_{A,N}(q,\cdot)$, so the explicit
  level-one truncated negative-moment estimate transfers uniformly to
  $F_{A,N}(q,\cdot)^{-\gamma N}$ for every $\gamma\in(0,1/2)$.
  Absorbing the fixed prefactor and the coefficient
  $c_{A,R_0}^{-\gamma N}$ into one exponential rate gives a real $b$ such
  that the moment is at most $e^{bN}$ for every $N\geq1$.
\end{lemma}
\begin{proof}\leanok\end{proof}

\begin{proposition}[Stationary negative-moment drift]
  \label{prop:gaussian-stationary-negative-moment-drift}
  \lean{AbsorptionCutoff.exists_eventually_integrable_neg_rpow_Fmap_foster,
    AbsorptionCutoff.exists_eventually_integrable_neg_rpow_Kchain_foster,
    AbsorptionCutoff.integrable_neg_rpow_Kchain_pow_of_foster,
    AbsorptionCutoff.exists_uniform_integral_neg_rpow_Kchain_pow_le_of_foster,
    AbsorptionCutoff.exists_uniform_integral_neg_rpow_cesaroMeasure_le_of_foster,
    AbsorptionCutoff.exists_eventually_uniform_integral_neg_rpow_Kchain_pow_and_cesaroMeasure_le,
    AbsorptionCutoff.negativeMomentTruncation,
    AbsorptionCutoff.negativeMomentTruncation_apply,
    AbsorptionCutoff.negativeMomentTruncation_nonneg,
    AbsorptionCutoff.negativeMomentTruncation_le,
    AbsorptionCutoff.negativeMomentTruncation_mono,
    AbsorptionCutoff.tendsto_negativeMomentTruncation,
    AbsorptionCutoff.integrable_neg_rpow_of_tendsto_of_uniform_integral,
    AbsorptionCutoff.exists_invariant_Kchain_integrable_neg_rpow_of_uniform_cesaro,
    AbsorptionCutoff.measureReal_Ioc_le_rpow_mul_integral_neg_rpow,
    AbsorptionCutoff.exists_pos_lt_and_add_mul_log_neg,
    AbsorptionCutoff.exists_exponential_envelope_invariant_moment_max,
    AbsorptionCutoff.exists_eventually_invariant_Kchain_integrable_neg_rpow,
    AbsorptionCutoff.exists_eventually_invariant_Kchain_small_ball_le_exp,
    AbsorptionCutoff.invariant_measureReal_le_add_of_ae_kernel_pow_measureReal_le,
    AbsorptionCutoff.add_mul_exp_neg_le_mul_add_exp_neg_min,
    AbsorptionCutoff.exists_eventually_invariant_Kchain_stable_interval_le_exp_add,
    AbsorptionCutoff.exists_eventually_invariant_Kchain_stable_interval_le_exp,
    AbsorptionCutoff.exists_eventually_invariant_Kchain_stable_interval_le_inv_nat,
    AbsorptionCutoff.integral_sq_Fmap_sub_V_le,
    AbsorptionCutoff.integral_sq_Fmap_sub_fixed_eq,
    AbsorptionCutoff.invariant_integral_sq_sub_fixed_eq_integral_integral_Fmap,
    AbsorptionCutoff.invariant_integral_sq_sub_fixed_eq_integral_add_noise,
    AbsorptionCutoff.invariant_integral_sq_sub_fixed_le_integral_sq_V_add,
    AbsorptionCutoff.abs_V_sub_fixed_le_mul_abs_sub_fixed_of_deriv_le,
    AbsorptionCutoff.integral_sq_V_sub_fixed_le_sq_mul_integral_add_measureReal,
    AbsorptionCutoff.exists_eventually_invariant_Kchain_integral_sq_sub_fixed_le_inv_nat}
  \leanok
  \uses{lem:gaussian-uniform-small-tanh-ratio,
    lem:gaussian-outside-coordinate-lower}
  For the exponent and rates selected above, and every sufficiently large
  positive dimension, the Lyapunov function
  $V_N(q)=q^{-\gamma N}$ satisfies on $(0,1]$
  \[
    K_{A,N}V_N(q)
    \leq e^{-\kappa N}V_N(q)+e^{bN}.
  \]
  The small-radius case factors out $q^{-\gamma N}$ and uses the ratio
  contraction; the complementary case uses the away-from-zero bound.
\end{proposition}
\begin{proof}\leanok\end{proof}

\begin{theorem}[Power singularity in fixed dimension]
  \label{thm:nd-power-singularity}
  \lean{AbsorptionCutoff.invariant_powerSingularity}
  \leanok
  \uses{lem:nd-gaussian-renewal,prop:nd-forcing-admissibility}
  For every origin-free invariant vector law in the supercritical regime, the
  Cram\'er exponent is the radial singularity exponent and the coefficient is
  strictly positive.  After multiplication by $s^{-\beta_{A,N}}$, directional
  small-ball masses converge on every measurable angular continuity set to the
  corresponding normalized surface mass times this coefficient; the full-ball
  tail converges to the coefficient itself.
\end{theorem}
\begin{proof}\leanok\end{proof}

\begin{definition}[Reduction kernels]
  \label{def:reduction-kernels}
  \lean{AbsorptionCutoff.radiusSq, AbsorptionCutoff.radiusSq_eq_zero_iff,
    AbsorptionCutoff.map_radiusSq_apply_singleton_zero, AbsorptionCutoff.Jmap,
    AbsorptionCutoff.Jkernel, AbsorptionCutoff.weightVar, AbsorptionCutoff.gaussianMat,
    AbsorptionCutoff.Pstep, AbsorptionCutoff.Pkernel}
  \leanok
  \uses{def:gaussianVec}
  The normalized squared radius $r_N(x)=N^{-1}\sum_i x_i^2$
  (\texttt{radiusSq}); in positive dimension its zero fiber is exactly the
  zero vector, so radius pushforward preserves the atom at the origin. The
  reconstruction map $(\tanh(A\sqrt q\,g_i))_i$ (\texttt{Jmap})
  and the Markov \emph{reconstruction kernel} $J_{A,N}(q,\cdot)$, the law of
  $(\tanh(A\sqrt q\,G_i))_i$ with $G\sim\mathcal G_N$ (\texttt{Jkernel}); and the vector
  chain of \texttt{eq:unrounded-chain}: the weight variance $A^2/N$ (\texttt{weightVar}),
  the $N\times N$ i.i.d.\ Gaussian weight law (\texttt{gaussianMat}), the step map
  $\tanh(\mathsf Wx)$ (\texttt{Pstep}), and the Markov \emph{unrounded vector kernel}
  $P_{A,N}(x,\cdot)$, the law of $\tanh(\mathsf Wx)$ (\texttt{Pkernel}).
\end{definition}

\begin{lemma}[Radius factorization of the reconstruction kernel]
  \label{lem:radius-factorization}
  \lean{AbsorptionCutoff.radiusSq_Jmap, AbsorptionCutoff.Jkernel_apply, AbsorptionCutoff.Kchain_apply,
    AbsorptionCutoff.Pkernel_apply, AbsorptionCutoff.radiusSq_map_Jkernel}
  \leanok
  \uses{def:reduction-kernels, def:Kchain, def:Fmap}
  For every $q\geq0$,
  \[
    (r_N)_\# J_{A,N}(q,\cdot)=K_{A,N}(q,\cdot).
  \]
\end{lemma}
\begin{proof}[Proof outline]
  \leanok
  \uses{def:reduction-kernels, def:Kchain}
  The pointwise identity $r_N(J_{A,N}(q,g))=F_{A,N}(q,g)$ intertwines the two
  Gaussian pushforward formulas.
\end{proof}

\begin{lemma}[Finite sums of independent real Gaussians]
  \label{lem:gaussian-sum}
  \lean{AbsorptionCutoff.map_finsetSum_gaussianReal, AbsorptionCutoff.map_sum_mul_gaussianPi}
  \leanok
  If $(X_i)$ are independent and $X_i\sim\mathcal N(m_i,v_i)$, then
  \[
    \sum_{i\in s}X_i
      \sim\mathcal N\!\left(\sum_{i\in s}m_i,\sum_{i\in s}v_i\right).
  \]
  In particular, under $\bigotimes_i\mathcal N(0,v)$ on $\mathbb R^N$,
  \[
    \sum_j w_jx_j\sim\mathcal N\!\left(0,v\sum_jx_j^2\right).
  \]
\end{lemma}
\begin{proof}[Proof outline]
  \leanok
  Induct over the finite index set using the two-variable Gaussian addition
  law; the second assertion is the corresponding linear-form specialization.
\end{proof}

\begin{lemma}[Gaussian isotropy / vector kernel factorization]
  \label{lem:gaussian-isotropy}
  \lean{AbsorptionCutoff.map_rowMap_gaussianMat, AbsorptionCutoff.map_tanh_scale_gaussian,
    AbsorptionCutoff.Jkernel_apply_pi, AbsorptionCutoff.radiusSq_nonneg,
    AbsorptionCutoff.Pkernel_eq_Jkernel_radiusSq}
  \leanok
  \uses{def:reduction-kernels, lem:gaussian-sum, lem:radius-factorization}
  The unrounded vector kernel depends on $x$ only through its normalized
  squared radius:
  \[
    P_{A,N}(x,\cdot)=J_{A,N}(r_N(x),\cdot).
  \]
\end{lemma}
\begin{proof}[Proof outline]
  \leanok
  \uses{lem:gaussian-sum, lem:radius-factorization}
  By \Cref{lem:gaussian-sum}, the rows of $\mathsf Wx$ are independent
  $\mathcal N(0,A^2r_N(x))$ variables. Apply $\tanh$ coordinatewise and use
  Gaussian scaling to identify the resulting product law with $J_{A,N}$.
\end{proof}

\begin{lemma}[Vector and radius laws at arbitrary times]
  \label{lem:gaussian-vector-radius-laws}
  \lean{AbsorptionCutoff.radiusSq_map_Pkernel_pow,
    AbsorptionCutoff.Pkernel_pow_succ_eq_Jkernel_comp_Kchain_pow}
  \leanok
  \uses{lem:gaussian-isotropy, lem:radius-factorization}
  For every $t\geq0$,
  \[
    (r_N)_\#P_{A,N}^t(x,\cdot)=K_{A,N}^t(r_N(x),\cdot),
    \qquad
    P_{A,N}^{t+1}(x,\cdot)=K_{A,N}^{t}(r_N(x),\cdot)J_{A,N}.
  \]
\end{lemma}
\begin{proof}[Proof outline]
  \leanok
  \uses{lem:gaussian-isotropy, lem:radius-factorization}
  Induct the radius identity from the one-step intertwining. For the second
  identity, expose the final vector step, factor it through the radius
  pushforward, and apply the first identity.
\end{proof}

\begin{lemma}[Invariant-law transfer]
  \label{lem:gaussian-invariant-transfer}
  \lean{AbsorptionCutoff.invariant_Pkernel_of_invariant_Kchain,
    AbsorptionCutoff.eq_Jkernel_comp_map_radiusSq_of_invariant_Pkernel,
    AbsorptionCutoff.invariant_Kchain_map_radiusSq_of_invariant_Pkernel}
  \leanok
  \uses{lem:gaussian-vector-radius-laws}
  If $\nu K_{A,N}=\nu$, then $\pi:=\nu J_{A,N}$ satisfies
  $\pi P_{A,N}=\pi$. Conversely, if $\pi P_{A,N}=\pi$, then
  \[
    \pi=((r_N)_\#\pi)J_{A,N},
    \qquad
    ((r_N)_\#\pi)K_{A,N}=(r_N)_\#\pi.
  \]
\end{lemma}
\begin{proof}[Proof outline]
  \leanok
  \uses{lem:gaussian-vector-radius-laws}
  Apply the two identities of \Cref{lem:gaussian-vector-radius-laws} and the
  assumed stationarity.
\end{proof}

\begin{lemma}[Weak stationary equation for the Gaussian vector chain]
  \label{lem:gaussian-vector-weak-stationary-equation}
  \lean{AbsorptionCutoff.integral_Pkernel, AbsorptionCutoff.vector_stationary_equation}
  \leanok
  \uses{def:reduction-kernels}
  If $\pi$ is invariant for $P_{A,N}$, then every bounded measurable test function
  $\varphi$ satisfies
  \[
    \int \varphi(x)\,\pi(dx)
    =
    \int\!\!\int \varphi\bigl(\operatorname{Pstep}_N(x,\mathsf W)\bigr)\,
      \operatorname{gaussianMat}_{A,N}(d\mathsf W)\,\pi(dx).
  \]
\end{lemma}
\begin{proof}[Proof outline]
  \leanok
  \uses{def:reduction-kernels}
  Use the pushforward formula for $P_{A,N}(x,\cdot)$, invariance, and the
  kernel--measure Fubini identity.
\end{proof}

\begin{lemma}[Gaussian vector--radius total-variation bracket]
  \label{lem:gaussian-tv-bracket}
  \lean{AbsorptionCutoff.tvDist_Kchain_pow_le_Pkernel_pow,
    AbsorptionCutoff.tvDist_Pkernel_pow_succ_le_Kchain_pow}
  \leanok
  \uses{lem:gaussian-vector-radius-laws, lem:gaussian-invariant-transfer}
  For $\pi=\nu J_{A,N}$ with $\nu$ invariant for $K_{A,N}$, the scalar radius
  observation gives the lower total-variation bound, while contraction through the
  reconstruction kernel $J_{A,N}$ gives the upper bound at positive times:
  \[
    \|K_{A,N}^t(r_Nx,\cdot)-\nu\|_{\mathrm{TV}}
    \le \|P_{A,N}^t(x,\cdot)-\pi\|_{\mathrm{TV}}
    \le \|K_{A,N}^{t-1}(r_Nx,\cdot)-\nu\|_{\mathrm{TV}}.
  \]
\end{lemma}
\begin{proof}\leanok\uses{lem:gaussian-vector-radius-laws, lem:gaussian-invariant-transfer}\end{proof}

\begin{lemma}[Interior transition density and invariant-law uniqueness]
  \label{lem:gaussian-one-coordinate-density}
  \lean{AbsorptionCutoff.gaussianOneCoordinateScore,
    AbsorptionCutoff.measurable_gaussianOneCoordinateScore,
    AbsorptionCutoff.gaussianOneCoordinateScore_map,
    AbsorptionCutoff.hasDerivAt_gaussianOneCoordinateDensityFormula_q,
    AbsorptionCutoff.integrable_gaussianOneCoordinateScore_map,
    AbsorptionCutoff.integral_gaussianOneCoordinateScore_map_eq_zero,
    AbsorptionCutoff.integral_fourth_gaussian,
    AbsorptionCutoff.integral_sq_gaussianOneCoordinateScore_map,
    AbsorptionCutoff.integral_pow_four_gaussianReal,
    AbsorptionCutoff.integrable_pow_four_gaussianReal,
    AbsorptionCutoff.integrable_sq_gaussianOneCoordinateScore_map_and_integral_eq,
    AbsorptionCutoff.gaussianCoordinateScoreSum,
    AbsorptionCutoff.integrable_gaussianOneCoordinateScore_map_eval,
    AbsorptionCutoff.integral_gaussianOneCoordinateScore_map_eval_eq_zero,
    AbsorptionCutoff.integrable_gaussianCoordinateScoreSum,
    AbsorptionCutoff.integral_gaussianCoordinateScoreSum_eq_zero,
    AbsorptionCutoff.integrable_sq_gaussianOneCoordinateScore_map_eval,
    AbsorptionCutoff.memLp_two_gaussianOneCoordinateScore_map,
    AbsorptionCutoff.variance_gaussianOneCoordinateScore_map,
    AbsorptionCutoff.integrable_sq_gaussianCoordinateScoreSum,
    AbsorptionCutoff.integral_sq_gaussianCoordinateScoreSum,
    AbsorptionCutoff.integrable_abs_gaussianCoordinateScoreSum,
    AbsorptionCutoff.integral_abs_gaussianCoordinateScoreSum_le,
    AbsorptionCutoff.gaussianProductDensity,
    AbsorptionCutoff.gaussianProductScore,
    AbsorptionCutoff.measurable_gaussianProductDensity,
    AbsorptionCutoff.measurable_gaussianProductScore,
    AbsorptionCutoff.measurable_gaussianProductScore_mul_density,
    AbsorptionCutoff.hasDerivAt_gaussianOneCoordinateDensity_q,
    AbsorptionCutoff.hasDerivAt_gaussianProductDensity_q,
    AbsorptionCutoff.deriv_gaussianProductDensity_q,
    AbsorptionCutoff.continuousAt_gaussianProductScore_mul_density_q,
    AbsorptionCutoff.gaussianProductObservationMap,
    AbsorptionCutoff.gaussianProductObservationMeasure,
    AbsorptionCutoff.measurable_gaussianProductObservationMap,
    AbsorptionCutoff.map_gaussianProductObservationMap_eq,
    AbsorptionCutoff.gaussianProductScore_observationMap,
    AbsorptionCutoff.integral_abs_gaussianProductScore_eq,
    AbsorptionCutoff.integral_abs_gaussianProductScore_le,
    AbsorptionCutoff.integrable_gaussianOneCoordinateDensity,
    AbsorptionCutoff.integrable_gaussianProductDensity,
    AbsorptionCutoff.gaussianProductObservationMeasure_eq_withDensity,
    AbsorptionCutoff.gaussianProductDensity_nonneg,
    AbsorptionCutoff.gaussianProductObservationMeasure_apply_toReal,
    AbsorptionCutoff.integrable_gaussianProductScore_observationMeasure,
    AbsorptionCutoff.integrable_gaussianProductScore_mul_density,
    AbsorptionCutoff.intervalIntegral_gaussianProductScore_mul_density_eq_sub,
    AbsorptionCutoff.abs_gaussianProductDensity_sub_le_integral_abs,
    AbsorptionCutoff.integral_abs_gaussianProductScore_mul_density_eq,
    AbsorptionCutoff.integral_abs_gaussianProductScore_mul_density_le,
    AbsorptionCutoff.integral_abs_gaussianProductScore_mul_density_le_uIcc,
    AbsorptionCutoff.integrable_gaussianProductScore_mul_density_prod,
    AbsorptionCutoff.integral_abs_gaussianProductDensity_sub_le,
    AbsorptionCutoff.tvDist_gaussianProductObservationMeasure_le,
    AbsorptionCutoff.gaussianProductAverage,
    AbsorptionCutoff.measurable_gaussianProductAverage,
    AbsorptionCutoff.gaussianProductAverage_observationMap,
    AbsorptionCutoff.map_gaussianProductAverage_observationMeasure_eq_Kchain,
    AbsorptionCutoff.tvDist_Kchain_apply_le_of_pos,
    AbsorptionCutoff.tanh_sq_sqrt_mul_mono,
    AbsorptionCutoff.monotoneOn_Fmap,
    AbsorptionCutoff.integral_abs_Fmap_sub_eq_abs_V_sub,
    AbsorptionCutoff.synchronousFmap,
    AbsorptionCutoff.measurable_synchronousFmap,
    AbsorptionCutoff.synchronousKchain,
    AbsorptionCutoff.synchronousKchain_apply,
    AbsorptionCutoff.synchronousKchain_map_fst,
    AbsorptionCutoff.synchronousKchain_map_snd,
    AbsorptionCutoff.integral_abs_fst_sub_snd_synchronousKchain,
    AbsorptionCutoff.synchronousKchain_apply_prod_Icc_compl,
    AbsorptionCutoff.markovPathMeasure_ae_eval_succ_mem_synchronousKchain_prod_Icc,
    AbsorptionCutoff.markovPathMeasure_dirac_ae_eval_mem_synchronousKchain_prod_Icc,
    AbsorptionCutoff.condExp_abs_fst_sub_snd_eval_succ_eq_abs_V_sub,
    AbsorptionCutoff.markovPathMeasure_map_eval_synchronousKchain_map_fst,
    AbsorptionCutoff.markovPathMeasure_map_eval_synchronousKchain_map_snd,
    AbsorptionCutoff.markovPathMeasure_measureReal_eval_not_mem_prod_le,
    AbsorptionCutoff.markovPathMeasure_measureReal_eval_not_mem_stableInterval_prod_le_add,
    AbsorptionCutoff.integral_sq_sub_V_Kchain_le,
    AbsorptionCutoff.integral_sq_sub_Kchain_eq,
    AbsorptionCutoff.integral_sq_sub_Kchain_le,
    AbsorptionCutoff.integral_sq_sub_V_Kchain_le_of_abs_V_sub_le,
    AbsorptionCutoff.integrable_sq_sub_of_ae_mem_Icc,
    AbsorptionCutoff.condExp_sq_eval_succ_sub_eq_integral_Kchain,
    AbsorptionCutoff.condExp_sq_eval_succ_sub_V_le_of_abs_V_sub_le,
    AbsorptionCutoff.integral_sq_eval_succ_sub_V_le_of_abs_V_sub_le,
    AbsorptionCutoff.integral_sq_eval_sub_V_iterate_le_of_forall_abs_V_sub_le,
    AbsorptionCutoff.stableIntervalExitTime,
    AbsorptionCutoff.stableIntervalExitTime_le_sentinel,
    AbsorptionCutoff.lt_stableIntervalExitTime_iff,
    AbsorptionCutoff.eval_mem_stableInterval_of_lt_stableIntervalExitTime,
    AbsorptionCutoff.abs_V_sub_le_mul_abs_sub_of_mem_stableInterval,
    AbsorptionCutoff.condExp_abs_fst_sub_snd_eval_succ_le_mul_add_indicator,
    AbsorptionCutoff.integrable_abs_fst_sub_snd_eval_synchronousKchain,
    AbsorptionCutoff.integral_abs_fst_sub_snd_eval_succ_le_mul_add_measureReal,
    AbsorptionCutoff.integral_abs_fst_sub_snd_eval_le_pow_mul_add,
    AbsorptionCutoff.integral_abs_fst_sub_snd_eval_le_pow_mul_add_of_scalar_localization,
    AbsorptionCutoff.markovPathMeasure_map_eval_synchronousKchain_map_fst_of_measure,
    AbsorptionCutoff.markovPathMeasure_map_eval_synchronousKchain_map_snd_of_measure,
    AbsorptionCutoff.markovPathMeasure_measureReal_eval_not_mem_prod_le_of_measure,
    AbsorptionCutoff.markovPathMeasure_measureReal_eval_not_mem_stableInterval_prod_le_add_of_measure,
    AbsorptionCutoff.markovPathMeasure_ae_eval_mem_synchronousKchain_prod_Icc_of_measure,
    AbsorptionCutoff.integrable_abs_fst_sub_snd_eval_synchronousKchain_of_measure,
    AbsorptionCutoff.condExp_abs_fst_sub_snd_eval_succ_eq_abs_V_sub_of_measure,
    AbsorptionCutoff.condExp_abs_fst_sub_snd_eval_succ_le_mul_add_indicator_of_measure,
    AbsorptionCutoff.integral_abs_fst_sub_snd_eval_succ_le_mul_add_measureReal_of_measure,
    AbsorptionCutoff.integral_abs_fst_sub_snd_eval_le_pow_mul_add_of_measure,
    AbsorptionCutoff.integral_abs_fst_sub_snd_eval_le_pow_mul_add_of_measure_of_scalar_localization,
    AbsorptionCutoff.Measure.prod_apply_compl_prod_Icc_eq_zero,
    AbsorptionCutoff.integral_abs_sub_le_sqrt_integral_sq,
    AbsorptionCutoff.integral_abs_fst_sub_snd_prod_le_add_sqrt_integral_sq,
    AbsorptionCutoff.integral_abs_fst_sub_snd_prod_le_add_sqrt_div_sqrt_nat,
    AbsorptionCutoff.integral_abs_fst_sub_snd_prod_le_sqrt_add_abs_add_sqrt,
    AbsorptionCutoff.integral_abs_fst_sub_snd_prod_le_sqrt_add_const_add_sqrt_div_sqrt_nat,
    AbsorptionCutoff.integral_abs_fst_sub_snd_map_eval_prod_le_sqrt_add_const_add_sqrt_div_sqrt_nat,
    AbsorptionCutoff.eventually_integral_abs_fst_sub_snd_blockStart_prod_le_inv_sqrt_nat,
    AbsorptionCutoff.exists_eventually_invariant_Kchain_family_integral_sq_sub_fixed_le_inv_nat,
    AbsorptionCutoff.exists_stationary_family_eventually_blockStart_prod_le_inv_sqrt_nat,
    AbsorptionCutoff.measureReal_abs_sub_gt_le_integral_sq_div,
    AbsorptionCutoff.measureReal_abs_sub_gt_le_div_sq_div_nat,
    AbsorptionCutoff.markovPathMeasure_map_eval_eq_of_invariant,
    AbsorptionCutoff.markovPathMeasure_measureReal_eval_mem_eq_of_invariant,
    AbsorptionCutoff.markovPathMeasure_measureReal_abs_eval_sub_gt_le_div_sq_div_nat,
    AbsorptionCutoff.eventually_forall_markovPathMeasure_measureReal_abs_eval_sub_gt_le_div_sq_div_nat,
    AbsorptionCutoff.markovPathMeasure_map_eval_of_map_eval,
    AbsorptionCutoff.markovPathMeasure_measureReal_eval_mem_of_map_eval,
    AbsorptionCutoff.markovPathMeasure_measureReal_eval_not_mem_stableInterval_le_exp,
    AbsorptionCutoff.eventually_forall_markovPathMeasure_measureReal_eval_not_mem_stableInterval_le_inv_nat,
    AbsorptionCutoff.eventually_forall_markovPathMeasure_shifted_measureReal_eval_not_mem_stableInterval_le_inv_nat,
    AbsorptionCutoff.markovPathMeasure_dirac_map_eval_apply_compl_Icc_eq_zero,
    AbsorptionCutoff.integral_abs_fst_sub_snd_eval_blockStart_prod_le_pow_mul_add_of_localization,
    AbsorptionCutoff.integral_abs_fst_sub_snd_eval_blockStart_prod_le_terminalContraction_pow_mul_add_of_localization,
    AbsorptionCutoff.measureReal_not_mem_stableInterval_le_four_mul_div_sq_div_nat_of_center,
    AbsorptionCutoff.markovPathMeasure_measureReal_eval_not_mem_stableInterval_le_four_mul_div_sq_div_nat_of_center,
    AbsorptionCutoff.eventually_forall_shifted_eval_not_mem_terminalInterval_le_of_center,
    AbsorptionCutoff.eventually_forall_integral_sq_eval_sub_V_iterate_le_inv_nat,
    AbsorptionCutoff.eventually_forall_integral_sq_eval_blockStart_add_le_inv_nat,
    AbsorptionCutoff.tendsto_mul_rpow_sub_terminalBlockLength_div_sqrt_div_terminalRadius_zero,
    AbsorptionCutoff.eventually_abs_V_iterate_terminalBlockStart_sub_fixed_le_half_terminalRadius,
    AbsorptionCutoff.forall_abs_V_iterate_sub_fixed_le_half_radius,
    AbsorptionCutoff.eventually_forall_abs_V_iterate_terminalBlockStart_add_sub_fixed_le_half_terminalRadius,
    AbsorptionCutoff.eventually_forall_shifted_eval_not_mem_terminalInterval_le_of_cutoff,
    AbsorptionCutoff.eventually_forall_stationary_eval_not_mem_terminalInterval_le_inv_nat,
    AbsorptionCutoff.eventually_integral_abs_fst_sub_snd_eval_blockStart_le_terminalContraction_pow_mul_add,
    AbsorptionCutoff.exists_eventually_integral_abs_fst_sub_snd_eval_terminalBlock_le_terminalContraction_pow_mul_add,
    AbsorptionCutoff.exists_eventually_terminalBlock_distance_le_terminalContraction_pow_mul_inv_sqrt_add,
    AbsorptionCutoff.exists_eventually_terminalBlock_distance_le_multiplier_pow_mul_inv_sqrt_add,
    AbsorptionCutoff.abs_V_sub_le_sq_mul_abs_sub_of_mem_Icc,
    AbsorptionCutoff.integral_sq_eval_succ_sub_V_iterate_le_global,
    AbsorptionCutoff.exists_nonneg_integral_sq_eval_sub_V_iterate_le_div_nat,
    AbsorptionCutoff.exists_eventually_markovPathMeasure_stable_entrance_exp_bound_of_tendsto,
    AbsorptionCutoff.markovPathMeasure_ae_eval_mem_Kchain_Icc_of_measure,
    AbsorptionCutoff.integral_sq_eval_succ_sub_V_le_of_abs_V_sub_le_of_measure,
    AbsorptionCutoff.condExp_sq_eval_succ_sub_V_iterate_le_on_lt_stableIntervalExitTime_of_measure,
    AbsorptionCutoff.condExp_indicator_sq_eval_succ_sub_V_iterate_le_of_measure,
    AbsorptionCutoff.stableIntervalKilledErrorMomentOfMeasure,
    AbsorptionCutoff.stableIntervalKilledErrorMomentOfMeasure_succ_le,
    AbsorptionCutoff.stableIntervalKilledErrorMomentOfMeasure_le,
    AbsorptionCutoff.markovPathMeasure_measureReal_abs_next_sub_V_gt_le_of_measure,
    AbsorptionCutoff.markovPathMeasure_measureReal_finiteHorizonKchainStepDeviationEvent_le_of_measure,
    AbsorptionCutoff.stableIntervalExitTime_eq_sentinel_of_initial_mem_of_step_deviation_le,
    AbsorptionCutoff.markovPathMeasure_measureReal_stableIntervalExitTime_le_of_measure,
    AbsorptionCutoff.integral_sq_eval_sub_V_iterate_le_killed_add_exit_of_measure,
    AbsorptionCutoff.integral_sq_eval_sub_V_iterate_le_initial_add_inv_add_exp_of_measure,
    AbsorptionCutoff.integral_sq_eval_add_sub_V_iterate_le_initial_add_inv_add_exp,
    AbsorptionCutoff.exists_uniform_markovPathMeasure_stable_entrance_exp_bound_with_margin,
    AbsorptionCutoff.exists_eventually_markovPathMeasure_stable_entrance_exp_bound_with_margin_of_tendsto,
    AbsorptionCutoff.exists_eventually_forall_integral_sq_eval_sub_V_iterate_le_inv_nat_of_tendsto,
    AbsorptionCutoff.exists_eventually_forall_shifted_eval_not_mem_terminalInterval_le_of_tendsto,
    AbsorptionCutoff.exists_eventually_integral_abs_fst_sub_snd_terminalBlockStart_prod_le_rpow_div_sqrt_of_tendsto,
    AbsorptionCutoff.exists_eventually_terminalBlock_distance_le_contraction_pow_mul_rpow_add_of_tendsto,
    AbsorptionCutoff.exists_eventually_terminalBlock_distance_le_multiplier_pow_mul_rpow_add_of_tendsto,
    AbsorptionCutoff.pow_mul_rpow_sub_natCast,
    AbsorptionCutoff.tendsto_pow_supercriticalTerminalBlockLength_zero,
    AbsorptionCutoff.tendsto_sqrt_nat_mul_terminalLocalizationError_zero,
    AbsorptionCutoff.tendsto_inv_nat_div_supercriticalTerminalRadius_sq_zero,
    AbsorptionCutoff.tvDist_kernel_comp_map_fst_snd_le_integral_add_measureReal,
    AbsorptionCutoff.tvDist_Kchain_comp_map_fst_snd_le_integral_add_stableInterval,
    AbsorptionCutoff.tvDist_Kchain_comp_map_eval_fst_snd_le_integral_add_stableInterval,
    AbsorptionCutoff.exists_eventually_shifted_eval_terminalBlockLength_not_mem_terminalInterval_le_of_tendsto,
    AbsorptionCutoff.exists_eventually_terminalBlock_endpoint_not_mem_prod_le_of_tendsto,
    AbsorptionCutoff.exists_eventually_sqrt_nat_mul_terminalBlock_distance_le_rpow_add,
    AbsorptionCutoff.exists_eventually_terminalBlock_smoothed_tvDist_le_rpow_add_of_tendsto,
    AbsorptionCutoff.eventually_terminalBlock_smoothed_fst_eq_cutoff_marginal,
    AbsorptionCutoff.terminalBlock_smoothed_snd_eq_of_invariant,
    AbsorptionCutoff.exists_eventually_tvDist_cutoff_marginal_invariant_le_rpow_add_of_tendsto,
    AbsorptionCutoff.exists_stationary_family_eventually_tvDist_cutoff_marginal_le_rpow_add_of_tendsto,
    AbsorptionCutoff.tvDist_Pkernel_pow_succ_le_markovPathMeasure_radius_marginal,
    AbsorptionCutoff.exists_reconstructed_invariant_vector_family_eventually_tvDist_cutoff_succ_le_rpow_add,
    AbsorptionCutoff.tendsto_normalized_V_orbit_sub_fixed_of_tendsto,
    AbsorptionCutoff.exists_eventually_rpow_le_sqrt_nat_mul_abs_V_iterate_integerCutoffTime_sub_fixed,
    AbsorptionCutoff.one_sub_four_mul_add_div_sq_le_tvDist_of_integral_sq,
    AbsorptionCutoff.eventually_one_sub_four_mul_add_div_sq_rpow_le_tvDist_of_inv_nat_moments,
    AbsorptionCutoff.exists_stationary_family_eventually_one_sub_mul_rpow_le_tvDist_cutoff_marginal,
    AbsorptionCutoff.exists_reconstructed_invariant_vector_family_eventually_one_sub_mul_rpow_le_tvDist_cutoff,
    AbsorptionCutoff.exists_forall_eventually_supercriticalIntegerCutoffTime_le_mul_log_nat,
    AbsorptionCutoff.forall_eventually_supercriticalIntegerCutoffTime_le_nat,
    AbsorptionCutoff.exists_forall_eventually_forall_integral_sq_eval_sub_V_iterate_le_inv_nat_cutoff,
    AbsorptionCutoff.exists_forall_eventually_rpow_le_sqrt_nat_mul_abs_V_iterate_integerCutoffTime_sub_fixed,
    AbsorptionCutoff.exists_forall_eventually_integrable_integral_sq_cutoffTime_sub_V_iterate_le_inv_nat,
    AbsorptionCutoff.exists_stationary_family_forall_eventually_one_sub_mul_rpow_le_tvDist_cutoff_marginal,
    AbsorptionCutoff.exists_reconstructed_invariant_vector_family_forall_eventually_one_sub_mul_rpow_le_tvDist_cutoff,
    AbsorptionCutoff.exists_forall_eventually_abs_V_iterate_sub_fixed_le_mul_pow_of_tendsto,
    AbsorptionCutoff.exists_forall_eventually_abs_V_iterate_terminalBlockStart_sub_fixed_le_mul_rpow_div_sqrt,
    AbsorptionCutoff.exists_forall_eventually_integrable_integral_sq_terminalBlockStart_sub_V_iterate_le_inv_nat,
    AbsorptionCutoff.exists_forall_eventually_integral_abs_fst_sub_snd_terminalBlockStart_prod_le_rpow_div_sqrt,
    AbsorptionCutoff.exists_forall_eventually_shifted_eval_not_mem_terminalInterval_le,
    AbsorptionCutoff.exists_forall_eventually_terminalBlock_distance_le_contraction_pow_mul_rpow_add,
    AbsorptionCutoff.exists_forall_eventually_terminalBlock_distance_le_multiplier_pow_mul_rpow_add,
    AbsorptionCutoff.exists_forall_eventually_sqrt_nat_mul_terminalBlock_distance_le_rpow_add,
    AbsorptionCutoff.exists_forall_eventually_shifted_eval_terminalBlockLength_not_mem_terminalInterval_le,
    AbsorptionCutoff.exists_forall_eventually_terminalBlock_endpoint_not_mem_prod_le,
    AbsorptionCutoff.exists_forall_eventually_terminalBlock_smoothed_tvDist_le_rpow_add_of_tendsto,
    AbsorptionCutoff.exists_forall_eventually_tvDist_cutoff_marginal_invariant_le_rpow_add_of_tendsto,
    AbsorptionCutoff.exists_stationary_family_forall_eventually_tvDist_cutoff_marginal_le_rpow_add_of_tendsto,
    AbsorptionCutoff.exists_reconstructed_invariant_vector_family_forall_eventually_tvDist_cutoff_succ_le_rpow_add,
    AbsorptionCutoff.eventually_supercriticalIntegerCutoffTime_sub_one_add_one_eq,
    AbsorptionCutoff.exists_reconstructed_invariant_vector_family_forall_eventually_tvDist_cutoff_le_rpow_add,
    AbsorptionCutoff.eventually_eq_of_invariant_Kchain_family_apply_singleton_zero,
    AbsorptionCutoff.exists_reconstructed_invariant_vector_family_forall_eventually_two_sided_tvDist_cutoff_profile,
    AbsorptionCutoff.exists_reconstructed_invariant_vector_family_forall_liminf_limsup_tvDist_cutoff_profile,
    AbsorptionCutoff.exists_reconstructed_invariant_vector_family_tendsto_outer_tvDist_cutoff_profile,
    AbsorptionCutoff.exists_reconstructed_invariant_vector_family_hasCutoff,
    AbsorptionCutoff.radiusSq_le_one_of_forall_mem_Icc,
    AbsorptionCutoff.eventually_radiusSq_mem_Ioc_of_tendsto_of_forall_mem_Icc,
    AbsorptionCutoff.exists_reconstructed_invariant_vector_family_hasCutoff_of_forall_mem_Icc,
    AbsorptionCutoff.exists_stationary_family_forall_eventually_two_sided_tvDist_cutoff_profile,
    AbsorptionCutoff.exists_stationary_family_tendsto_outer_tvDist_cutoff_profile,
    AbsorptionCutoff.exists_stationary_family_hasCutoff,
    AbsorptionCutoff.exists_stationary_family_hasCutoff_of_forall_mem_Icc,
    AbsorptionCutoff.eventually_integral_abs_fst_sub_snd_eval_blockStart_prod_le_pow_mul_add_inv_nat,
    AbsorptionCutoff.eventually_integral_abs_fst_sub_snd_eval_blockStart_prod_le_pow_mul_inv_sqrt_add_inv_nat,
    AbsorptionCutoff.exists_stationary_family_eventually_terminal_distance_le_pow_mul_inv_sqrt_add_inv_nat,
    AbsorptionCutoff.supercriticalCutoffTime,
    AbsorptionCutoff.supercriticalCutoffTime_eq,
    AbsorptionCutoff.tendsto_supercriticalCutoffTime_atTop,
    AbsorptionCutoff.supercriticalIntegerCutoffTime,
    AbsorptionCutoff.supercriticalIntegerCutoffTime_eq,
    AbsorptionCutoff.rpow_supercriticalCutoffTime_eq_inv_sqrt_mul_abs_koenigsCoefficient,
    AbsorptionCutoff.supercriticalTerminalBlockLength,
    AbsorptionCutoff.supercriticalTerminalBlockLength_eq,
    AbsorptionCutoff.tendsto_supercriticalTerminalBlockLength_atTop,
    AbsorptionCutoff.tendsto_supercriticalTerminalBlockLength_div_log_nat_zero,
    AbsorptionCutoff.tendsto_supercriticalTerminalBlockLength_mul_terminalRadius_zero,
    AbsorptionCutoff.tendsto_supercriticalTerminalContraction_div_multiplier_pow_blockLength_one,
    AbsorptionCutoff.exists_abs_deriv_V_le_multiplier_add_mul_abs_sub_fixed,
    AbsorptionCutoff.exists_abs_V_sub_le_multiplier_add_mul_radius,
    AbsorptionCutoff.exists_eventually_supercriticalTerminal_contraction_data,
    AbsorptionCutoff.supercriticalTerminalBlockStart,
    AbsorptionCutoff.supercriticalTerminalBlockStart_eq,
    AbsorptionCutoff.supercriticalTerminalRadius,
    AbsorptionCutoff.supercriticalTerminalRadius_eq,
    AbsorptionCutoff.supercriticalTerminalRadius_pos,
    AbsorptionCutoff.tendsto_supercriticalTerminalRadius_zero,
    AbsorptionCutoff.supercriticalTerminalContraction,
    AbsorptionCutoff.supercriticalTerminalContraction_eq,
    AbsorptionCutoff.supercriticalTerminalContraction_nonneg,
    AbsorptionCutoff.tendsto_supercriticalTerminalContraction,
    AbsorptionCutoff.eventually_supercriticalTerminalContraction_lt_one,
    AbsorptionCutoff.tendsto_sqrt_nat_mul_inv_nat_div_supercriticalTerminalRadius_sq_zero,
    AbsorptionCutoff.tendsto_supercriticalIntegerCutoffTime_atTop,
    AbsorptionCutoff.tendsto_supercriticalCutoffTime_div_log_nat,
    AbsorptionCutoff.tendsto_supercriticalIntegerCutoffTime_div_log_nat,
    AbsorptionCutoff.eventually_supercriticalTerminalBlockLength_add_one_le_integerCutoffTime,
    AbsorptionCutoff.eventually_supercriticalTerminalBlockStart_add_length_eq_integerCutoffTime_sub_one,
    AbsorptionCutoff.eventually_pow_terminalBlockStart_le_cutoff_rpow,
    AbsorptionCutoff.tendsto_supercriticalTerminalBlockStart_div_log_nat,
    AbsorptionCutoff.tendsto_supercriticalTerminalBlockStart_atTop,
    AbsorptionCutoff.exists_eventually_abs_V_iterate_terminalBlockStart_sub_fixed_le_mul_pow,
    AbsorptionCutoff.exists_eventually_abs_V_iterate_terminalBlockStart_sub_fixed_le_mul_rpow_div_sqrt,
    AbsorptionCutoff.exists_eventually_supercriticalIntegerCutoffTime_le_mul_log_nat,
    AbsorptionCutoff.exists_eventually_supercriticalCutoff_block_horizon,
    AbsorptionCutoff.abs_V_eval_sub_V_iterate_le_of_lt_stableIntervalExitTime,
    AbsorptionCutoff.indicator_sq_succ_lt_stableIntervalExitTime_le,
    AbsorptionCutoff.measurableSet_eval_mem_stableInterval,
    AbsorptionCutoff.measurableSet_lt_stableIntervalExitTime,
    AbsorptionCutoff.measurableSet_stableIntervalExitTime_le,
    AbsorptionCutoff.condExp_sq_eval_succ_sub_V_iterate_le_on_lt_stableIntervalExitTime,
    AbsorptionCutoff.condExp_indicator_sq_eval_succ_sub_V_iterate_le,
    AbsorptionCutoff.stableIntervalKilledErrorMoment,
    AbsorptionCutoff.stableIntervalKilledErrorMoment_succ_le,
    AbsorptionCutoff.stableIntervalKilledErrorMoment_le,
    AbsorptionCutoff.stableIntervalExitTime_eq_sentinel_of_step_deviation_le,
    AbsorptionCutoff.markovPathMeasure_measureReal_stableIntervalExitTime_le,
    AbsorptionCutoff.integral_sq_eval_sub_V_iterate_le_killed_add_exit,
    AbsorptionCutoff.integral_sq_eval_sub_V_iterate_le_inv_add_exp,
    AbsorptionCutoff.eventually_integral_sq_eval_sub_V_iterate_le_inv_nat,
    AbsorptionCutoff.map_gaussianOneCoordinateInverse_gaussianWeight,
    AbsorptionCutoff.gaussianOneCoordinate_positiveBranch_weight,
    AbsorptionCutoff.map_gaussianOneCoordinateMap_restrict_Ioi,
    AbsorptionCutoff.map_gaussianOneCoordinateMap_restrict_Iio_eq_Ioi,
    AbsorptionCutoff.map_gaussianOneCoordinateMap_eq_withDensity,
    AbsorptionCutoff.gaussianTwoCoordinateDensity,
    AbsorptionCutoff.measurable_gaussianTwoCoordinateDensity,
    AbsorptionCutoff.conv_gaussianOneCoordinateMeasure_eq_withDensity,
    AbsorptionCutoff.gaussianTwoCoordinateDensity_nonneg,
    AbsorptionCutoff.gaussianTwoCoordinateDensity_pos,
    AbsorptionCutoff.gaussianCoordinateSumDensity,
    AbsorptionCutoff.gaussianCoordinateSumDensity_zero,
    AbsorptionCutoff.gaussianCoordinateSumDensity_succ,
    AbsorptionCutoff.gaussianCoordinateSumDensity_one,
    AbsorptionCutoff.measurable_gaussianCoordinateSumDensity,
    AbsorptionCutoff.gaussianCoordinateSumDensity_nonneg,
    AbsorptionCutoff.conv_gaussianCoordinateSumMeasure_succ_eq_withDensity,
    AbsorptionCutoff.gaussianCoordinateSumDensity_pos,
    AbsorptionCutoff.gaussianCoordinateSumLaw,
    AbsorptionCutoff.gaussianCoordinateSumLaw_zero,
    AbsorptionCutoff.gaussianCoordinateSumLaw_succ,
    AbsorptionCutoff.gaussianCoordinateSumLaw_eq_withDensity,
    AbsorptionCutoff.gaussianCoordinateSumMap,
    AbsorptionCutoff.measurable_gaussianCoordinateSumMap,
    AbsorptionCutoff.map_gaussianCoordinateSumMap_gaussianVec_succ,
    AbsorptionCutoff.map_gaussianCoordinateSumMap_gaussianVec,
    AbsorptionCutoff.gaussianAverageDensity,
    AbsorptionCutoff.measurable_gaussianAverageDensity,
    AbsorptionCutoff.Kchain_apply_eq_withDensity_gaussianAverageDensity,
    AbsorptionCutoff.Kchain_apply_absolutelyContinuous_volume,
    AbsorptionCutoff.volume_restrict_Ioo_absolutelyContinuous_Kchain_apply,
    AbsorptionCutoff.Kchain_apply_restrict_Ioo_absolutelyContinuous_volume_restrict_Ioo,
    AbsorptionCutoff.volume_restrict_Ioo_absolutelyContinuous_Kchain_apply_restrict_Ioo,
    AbsorptionCutoff.gaussianAverageDensity_pos,
    AbsorptionCutoff.Kchain_apply_pos_of_volume_pos,
    AbsorptionCutoff.Kchain_apply_Ioo_eq_one,
    AbsorptionCutoff.KchainInterior,
    AbsorptionCutoff.KchainInterior_apply,
    AbsorptionCutoff.KchainInterior_map_subtype_val,
    AbsorptionCutoff.map_subtype_val_apply_Ioo,
    AbsorptionCutoff.invariant_Kchain_map_subtype_val_of_invariant_KchainInterior,
    AbsorptionCutoff.map_comap_subtype_val_eq_of_apply_Ioo_eq_one,
    AbsorptionCutoff.isProbabilityMeasure_comap_subtype_val_of_apply_Ioo_eq_one,
    AbsorptionCutoff.invariant_KchainInterior_comap_subtype_val_of_invariant_Kchain,
    AbsorptionCutoff.invariant_pow,
    AbsorptionCutoff.irreducibleMeasure_absolutelyContinuous_invariant,
    AbsorptionCutoff.invariant_probability_unique,
    AbsorptionCutoff.KchainInteriorVolume,
    AbsorptionCutoff.KchainInteriorVolume_apply,
    AbsorptionCutoff.KchainInteriorVolume_ne_zero,
    AbsorptionCutoff.invariant_probability_unique_KchainInterior,
    AbsorptionCutoff.invariant_probability_unique_Kchain_of_apply_Ioo_eq_one,
    AbsorptionCutoff.invariant_Kchain_apply_Ioo_eq_one_of_Ioc_compl,
    AbsorptionCutoff.nonzeroPart_apply_Ioo_eq_one,
    AbsorptionCutoff.nonzeroPart_eq_of_invariant_Kchain}
  \leanok
  Let $A,q>0$. For a standard Gaussian $G$,
  \[
    (\tanh^2(A\sqrt q\,\cdot))_\#\mathcal N(0,1)
    =
    \mathrm{volume.withDensity}
      (\texttt{gaussianOneCoordinateDensity}\ A\ q).
  \]
  For $N>0$, the one-step radius kernel has density
  \[
    r\longmapsto
    N\,\texttt{gaussianCoordinateSumDensity}(A,q,N-1)(Nr),
  \]
  which is measurable and strictly positive on $(0,1)$. Hence the transition
  law and Lebesgue measure are mutually absolutely continuous on $(0,1)$, and
  the restricted kernel \texttt{KchainInterior} is irreducible with respect
  to the nonzero interior Lebesgue measure.

  The interior kernel has at most one invariant probability. Consequently,
  any two invariant probabilities for the full radius kernel that give mass
  one to $(0,1)$ coincide. More generally, any two invariant probabilities
  carried by $[0,1]$ and not concentrated at the absorbing origin have the
  same normalized nonzero component.
\end{lemma}
\begin{proof}[Proof outline]
  \leanok
  Apply the inverse-branch Jacobian formula on the positive half-line and use
  Gaussian symmetry for the negative branch. Recursive convolution gives the
  coordinate-sum densities and proves their strict positivity on the interior
  of the support. Scaling by $1/N$ yields the stated kernel density and mutual
  absolute continuity.

  Restriction to $(0,1)$ gives irreducibility in one step. A general
  measure-infimum argument shows that a kernel irreducible with respect to a
  nonzero measure has at most one invariant probability. Invariant measures
  are transported between the interior subtype and the full radius space by
  the inclusion map. Stationarity moves all nonzero mass into $(0,1)$, so the
  uniqueness result applies to normalized nonzero components.
\end{proof}

\begin{proposition}[Compactness selection and nonzero-component decomposition]
  \label{prop:gaussian-compactness-selection}
  \lean{AbsorptionCutoff.cesaroMeasure, AbsorptionCutoff.invariant_of_cesaro_weakLimit,
    AbsorptionCutoff.exists_invariant_of_cesaro_with_subseq,
    AbsorptionCutoff.exists_invariant_of_cesaro, AbsorptionCutoff.stationary_equation,
    AbsorptionCutoff.Kchain_zero, AbsorptionCutoff.Kchain_singleton_zero,
    AbsorptionCutoff.invariant_Kchain_apply_Icc_compl,
    AbsorptionCutoff.exists_invariant_Pkernel_of_cesaro,
    AbsorptionCutoff.measure_eq_atom_zero_add_restrict_compl,
    AbsorptionCutoff.invariant_restrict_compl_zero, AbsorptionCutoff.nonzeroPart,
    AbsorptionCutoff.nonzeroPart_eq_self_of_apply_singleton_zero,
    AbsorptionCutoff.inverseMomentTruncation,
    AbsorptionCutoff.inverseMomentTruncation_apply,
    AbsorptionCutoff.inverseMomentTruncation_zero,
    AbsorptionCutoff.inverseMomentTruncation_nonneg,
    AbsorptionCutoff.inverseMomentTruncation_le_one,
    AbsorptionCutoff.inverseMomentTruncation_le_scaled_inv,
    AbsorptionCutoff.apply_singleton_zero_of_tendsto_of_uniform_integral_inv,
    AbsorptionCutoff.Kchain_pow_apply_Ioo_eq_one,
    AbsorptionCutoff.cesaroPM_apply_Ioo_eq_one,
    AbsorptionCutoff.exists_invariant_Kchain_apply_singleton_zero_of_dimension,
    AbsorptionCutoff.exists_unique_invariant_probability_Kchain_of_apply_singleton_zero,
    AbsorptionCutoff.nonzeroPart_isProbabilityMeasure, AbsorptionCutoff.invariant_nonzeroPart,
    AbsorptionCutoff.nonzeroPart_Ioc_compl, AbsorptionCutoff.invariant_Pkernel_nonzeroPart,
    AbsorptionCutoff.invariant_Pkernel_nonzeroPart_unique,
    AbsorptionCutoff.invariant_Pkernel_nonzeroPart_unique_of_invariant_Pkernel,
    AbsorptionCutoff.invariant_probability_unique_Pkernel_of_apply_singleton_zero,
    AbsorptionCutoff.exists_unique_invariant_probability_Pkernel_of_apply_singleton_zero}
  \leanok
  \uses{def:Kchain, def:reduction-kernels, lem:gaussian-invariant-transfer,
    lem:gaussian-vector-weak-stationary-equation}
  Cesàro averages of $K_{A,N}$ started in $[0,1]$ are supported there and tight.
  They admit a weakly convergent subsequence whose limit $\nu$ is invariant
  and carried by $[0,1]$. Reconstruction through $J_{A,N}$ gives an invariant
  vector law.

  Every finite invariant radius measure splits as
  \[
    \nu=\nu\{0\}\delta_0+\nu|_{\{0\}^{c}},
  \]
  where the second summand is invariant. If $\nu\{0\}<1$, then
  \[
    \nu^+=(1-\nu\{0\})^{-1}\nu|_{\{0\}^{c}}
  \]
  is an invariant probability carried by $(0,1]$.

  If the approximating laws have uniformly bounded reciprocal moments, the
  weak limit satisfies $\nu\{0\}=0$. Under the explicit supercritical
  dimension criterion, there is therefore a unique origin-free invariant
  scalar probability and a unique origin-free invariant vector probability.
\end{proposition}
\begin{proof}[Proof outline]
  \leanok
  \uses{lem:gaussian-invariant-transfer, lem:gaussian-vector-weak-stationary-equation}
  Prokhorov compactness gives the weak limit; portmanteau preserves its support,
  and the Feller--Cesàro argument proves invariance. Absorption at zero permits
  cancellation of the atomic summand in the stationary equation.

  To exclude an atom at zero, use the bounded continuous truncations
  \[
    \tau_k(q)=\frac{1}{1+(k+1)|q|}
  \]
  which retain value one at the origin and, for $q>0$, satisfy
  $0\le\tau_k(q)\le(k+1)^{-1}q^{-1}$. Consequently, a weak limit of positive
  probability laws with uniformly bounded integrable reciprocal moments has
  no atom at zero. The interior irreducibility result gives uniqueness of the
  normalized nonzero component. Reconstruction transfers this uniqueness to
  the vector chain.
\end{proof}

\begin{theorem}[Supercritical dimension cutoff]
  \label{thm:gaussian-process-cutoff}
  \lean{AbsorptionCutoff.exists_stationary_family_hasCutoff_of_forall_mem_Icc,
    AbsorptionCutoff.exists_unique_nonzero_invariant_family_hasCutoff_of_forall_mem_Icc,
    AbsorptionCutoff.isCutoffWindow_supercriticalCutoffTime_one,
    AbsorptionCutoff.exists_unique_invariant_probability_Kchain_of_apply_singleton_zero}
  \leanok
  \uses{def:Kchain, def:cutoff, prop:gaussian-compactness-selection,
    lem:gaussian-vector-radius-laws}
  Let $A>1$, let $q_\ast\in(0,1)$ be the attracting nonzero fixed point of
  $V_A$, and let deterministic $x_N\in[-1,1]^N$ satisfy
  $N^{-1}\lVert x_N\rVert_2^2\to q_0\in(0,1]$, where $q_0\ne q_\ast$.
  Then the origin-free invariant scalar laws selected for $K_{A,N}$ have cutoff
  in the full sense of Definition~\ref{def:cutoff} at $t_N(q_0)$ with constant
  window one. In particular, the center diverges, the window is positive and
  little-o of the center, and the paper's two iterated total-variation limits
  hold. The interior irreducibility theorem
  identifies each selected law with the unique nonzero invariant law for all
  sufficiently large $N$.
\end{theorem}
\begin{proof}\leanok\end{proof}

\begin{corollary}[Vector supercritical cutoff]
  \label{cor:gaussian-vector-cutoff}
  \lean{AbsorptionCutoff.exists_reconstructed_invariant_vector_family_hasCutoff_of_forall_mem_Icc,
    AbsorptionCutoff.exists_unique_nonzero_invariant_vector_family_cutoff_of_forall_mem_Icc,
    AbsorptionCutoff.exists_unique_invariant_probability_Pkernel_of_apply_singleton_zero}
  \leanok
  \uses{thm:gaussian-process-cutoff, lem:gaussian-vector-radius-laws,
    lem:gaussian-invariant-transfer}
  Under the same assumptions, reconstructing the scalar stationary law
  through $J_{A,N}$ gives the unique origin-free invariant vector law.  The
  vector chain has the same full cutoff statement, center, and constant window
  one.
\end{corollary}
\begin{proof}\leanok\end{proof}

\chapter{Fixed-width cutoff}

This chapter records the fixed-dimension, vanishing-mesh cutoff route of paper
§3.  A positive-drift first-passage profile controls entrance to the rounding
scale; synchronous comparison and post-entrance absorption transfer that profile
to the rounded vector chain.

\begin{theorem}[Positive-drift first-passage profile]
  \label{thm:positive-drift-first-passage-profile}
  \lean{AbsorptionCutoff.tendsto_measureReal_firstPassageTime_gt,
    AbsorptionCutoff.tendsto_measureReal_firstPassageTime_gt_canonicalTime,
    AbsorptionCutoff.tendsto_measureReal_firstPassageTime_gt_postFloorTime}
  \leanok
  For an integrable iid increment process with positive mean and nonzero finite
  variance, the probability that the first-passage time above a diverging level
  exceeds its centered square-root-window observation time converges to the
  corresponding standard Gaussian distribution function.  The canonical-floor
  and post-floor forms allow any perturbation negligible on the square-root
  scale.
\end{theorem}
\begin{proof}\leanok\end{proof}

\begin{theorem}[Rounded fixed-width survival profile]
  \label{thm:fixed-width-rounded-survival}
  \lean{AbsorptionCutoff.tendsto_measureReal_fixedWidthRoundedSurvival_postFloorTime}
  \leanok
  \uses{thm:positive-drift-first-passage-profile}
  In the fixed-width subcritical regime, the rounded Gaussian recursion started
  from a nonzero vector survives through its post-floor cutoff observation time
  with probability converging to $\Phi(-a)$.  The proof combines the unrounded
  first-passage profile with synchronous coupling, a polylogarithmic entrance
  buffer, and uniform absorption after entrance.
\end{theorem}
\begin{proof}\leanok\uses{thm:positive-drift-first-passage-profile}\end{proof}

\begin{theorem}[Fixed-width total-variation profile]
  \label{thm:fixed-width-tv-profile}
  \lean{AbsorptionCutoff.tendsto_tvDist_roundedPkernel_postFloorTime,
    AbsorptionCutoff.tendsto_tvDist_roundedPkernel_fixedWidthMesh}
  \leanok
  \uses{thm:fixed-width-rounded-survival, lem:Hkernel-tv}
  Total variation from the absorbing origin equals the canonical rounded-chain
  survival probability.  Hence every diverging exponential mesh scale has
  Gaussian profile $\Phi(-a)$.  Equivalently, for every positive mesh sequence
  $\rho_r\to0$, the same limit holds at the manuscript's exact natural-floor
  time with
  \[
    L_{\rho_r}=\log\!\left(\frac{\lVert x_0\rVert_2}{\rho_r}\right).
  \]
\end{theorem}
\begin{proof}
  \leanok
  \uses{thm:fixed-width-rounded-survival, lem:Hkernel-tv}
\end{proof}

\begin{theorem}[Fixed-width cutoff and mixing times]
  \label{thm:fixed-width-cutoff}
  \lean{AbsorptionCutoff.hasCutoff_roundedPkernel_fixedWidth,
    AbsorptionCutoff.rounded_gaussian_nearest_cutoff_paper,
    AbsorptionCutoff.exists_eventually_fixedWidth_mixingTime_bounds}
  \leanok
  \uses{thm:fixed-width-tv-profile, def:cutoff, def:mixing}
  The rounded vector kernels have total-variation cutoff in the full sense of
  Definition~\ref{def:cutoff} at center $L/\mu$ and
  window $\sigma\mu^{-3/2}\sqrt L$.  For every fixed
  $\varepsilon\in(0,1)$, two constant Gaussian-profile offsets eventually
  bracket the mixing time, giving
  $t_{\mathrm{mix}}(\varepsilon)=L/\mu+O_\varepsilon(\sqrt L)$ with the stated
  variance prefactor.
\end{theorem}
\begin{proof}
  \leanok
  \uses{thm:fixed-width-tv-profile, def:cutoff, def:mixing}
\end{proof}

\chapter{The \texorpdfstring{$N$}{N}-dimensional supercritical stationary equation}

This chapter records paper §7.  In log-polar coordinates the stationary equation
becomes a Markov additive renewal problem; exponential tilting at the Cram\'er
root turns it into an ordinary renewal equation on the line, and the key renewal
theorem identifies the limit.

Mathlib has no renewal theory, so the linear-renewal input is proved in a
dedicated family of project modules. \emph{Deviation, deliberate and recorded:} the key
renewal theorem below is proved by approximation in the directly-Riemann-integrable
norm rather than through Feller's interval theorem, and therefore assumes the
kernel \emph{continuous} where the paper assumes only a.e.\ continuity.  Indicators
are not approximable in $\lVert\cdot\rVert_{\mathrm{DRI}}$ by continuous functions,
so the gap is real; it is harmless here because the chapter's forcing is
Gaussian-driven, hence atomless.  Continuity of the forcing is an explicit
obligation at each application.

\section{The Gaussian transfer moment and the Cram\'er exponent}

\begin{definition}[Gaussian transfer moment]
  \label{def:nd-gaussian-transfer-moment}
  \lean{AbsorptionCutoff.gaussianTransferMoment}
  \leanok
  $\mathcal M_{A,N}(\beta) = \mathbb E\lVert (A/\sqrt N)G\rVert_2^{-\beta}$, the
  spectral quantity to which isotropy collapses the transfer operator
  (\texttt{eq:nd-gaussian-negative-moment}).
\end{definition}

\begin{definition}[Log-radial increment]
  \label{def:nd-log-radial-increment}
  \lean{AbsorptionCutoff.logRadialIncrement}
  \leanok
  $Z = -\log\ell_{A,N}$, the logarithmic increment of the radial part.
\end{definition}

\begin{lemma}[Cram\'er exponent]
  \label{lem:nd-gaussian-cramer-exponent}
  \lean{AbsorptionCutoff.existsUnique_gaussianTransferMoment_eq_one,
    AbsorptionCutoff.cramerExponent, AbsorptionCutoff.cramerExponent_mem}
  \leanok
  \uses{def:nd-gaussian-transfer-moment, def:nd-log-radial-increment}
  For $A>A_c(N)$ there is a unique $\beta_{A,N}\in(0,N)$ with
  $r_{A,N}(\beta_{A,N})=1$, by strict convexity of $\beta\mapsto\log\mathcal
  M_{A,N}(\beta)$ together with $F(0)=0$, $F'(0)<0$ and $F\to\infty$ as
  $\beta\uparrow N$.
\end{lemma}
\begin{proof}\leanok\end{proof}

\section{The tilted increment law}

\begin{definition}[Tilted increment law]
  \label{def:nd-tilted-increment-law}
  \lean{AbsorptionCutoff.tiltedIncrementLaw}
  \leanok
  \uses{lem:nd-gaussian-cramer-exponent}
  $\widehat\mu_{A,N}$, the law of $Z$ after exponential tilting at
  $\beta_{A,N}$ (\texttt{eq:nd-tilted-increment-law}).
\end{definition}

\begin{lemma}[Properties of the tilted law]
  \label{lem:nd-tilted-increment-law-properties}
  \lean{AbsorptionCutoff.isProbabilityMeasure_tiltedIncrementLaw,
    AbsorptionCutoff.tiltedIncrementLaw_not_concentrated_on_lattice,
    AbsorptionCutoff.integral_id_tiltedIncrementLaw_pos,
    AbsorptionCutoff.memLp_id_two_tiltedIncrementLaw,
    AbsorptionCutoff.expTransform_convPow_tiltedIncrementLaw}
  \leanok
  \uses{def:nd-tilted-increment-law}
  $\widehat\mu_{A,N}$ is a nonlattice probability measure with drift
  $\widehat m_{A,N}\in(0,\infty)$ and a finite second moment, and
  $\mathbb E e^{-\beta_{A,N}S_n}=1$ for every $n$.  Nonlatticeness comes from
  atomlessness: the radial law has a density, so every fibre of $Z$ is null and
  every lattice, being countable, is null.
\end{lemma}
\begin{proof}\leanok\end{proof}

\section{The renewal mini-library}

\begin{definition}[Renewal measure and the d.R.i.\ norm]
  \label{def:nd-renewal-measure}
  \lean{AbsorptionCutoff.Renewal.convPow, AbsorptionCutoff.Renewal.renewalMeasure,
    AbsorptionCutoff.Renewal.driNorm, AbsorptionCutoff.Renewal.expTransform}
  \leanok
  $U=\sum_{n\ge0}\mu^{*n}$, and the directly-Riemann-integrable norm
  $\lVert z\rVert_{\mathrm{DRI}}=\sum_{k\in\mathbb Z}\sup_{[k,k+1]}|z|$ of
  \texttt{eq:nd-dri-definition}.
\end{definition}

\begin{theorem}[Key renewal theorem]
  \label{thm:nd-key-renewal}
  \lean{AbsorptionCutoff.Renewal.tendsto_tsum_integral_comp_sub_of_driNorm,
    AbsorptionCutoff.Renewal.tendsto_tsum_integral_comp_sub_of_driNorm_real}
  \leanok
  \uses{def:nd-renewal-measure}
  For a nonlattice probability measure with a second moment, positive drift and
  a finite exponential moment, and for continuous $z$ with
  $\lVert z\rVert_{\mathrm{DRI}}<\infty$,
  \[
    \sum_n\int z(y-s)\,\mu^{*n}(\mathrm ds)\longrightarrow
    \frac{1}{\widehat m}\int_{\mathbb R} z ,\qquad y\to\infty .
  \]
  Proved by Abel-regularized Fourier inversion against bandlimited kernels,
  followed by approximation in $\lVert\cdot\rVert_{\mathrm{DRI}}$.
\end{theorem}
\begin{proof}\leanok\end{proof}

\begin{theorem}[Key renewal theorem at $\widehat\mu_{A,N}$]
  \label{thm:nd-key-renewal-tilted}
  \lean{AbsorptionCutoff.tendsto_tsum_integral_comp_sub_tiltedIncrementLaw,
    AbsorptionCutoff.tendsto_tsum_integral_comp_sub_tiltedIncrementLaw_real}
  \leanok
  \uses{thm:nd-key-renewal, lem:nd-tilted-increment-law-properties}
  The instantiation of the previous theorem at the tilted increment law: all
  four hypotheses are supplied by
  \cref{lem:nd-tilted-increment-law-properties}.
\end{theorem}
\begin{proof}\leanok\end{proof}

\section{Solving the renewal equation}

\begin{lemma}[Iterating the renewal equation]
  \label{lem:nd-renewal-iteration}
  \lean{AbsorptionCutoff.Renewal.eq_integral_convPow_add_sum_of_renewalEquation}
  \leanok
  \uses{def:nd-renewal-measure}
  If $h_y=\int h_{y-z}\,\mu(\mathrm dz)+\psi_y$ then
  $h_y=\mathbb E h_{y-S_n}+\sum_{k<n}\mathbb E\psi_{y-S_k}$ for every $n$.
\end{lemma}
\begin{proof}\leanok\end{proof}

\begin{lemma}[The terminal term vanishes]
  \label{lem:nd-renewal-terminal}
  \lean{AbsorptionCutoff.Renewal.tendsto_convPow_Icc_zero,
    AbsorptionCutoff.Renewal.lintegral_enorm_comp_sub_le,
    AbsorptionCutoff.Renewal.tendsto_integral_comp_sub_convPow_zero,
    AbsorptionCutoff.Renewal.eq_tsum_integral_comp_sub_of_renewalEquation}
  \leanok
  \uses{lem:nd-renewal-iteration, lem:nd-tilted-increment-law-properties}
  Under the left-boundary condition, local boundedness, and right-tail
  minimality, $\mathbb E h_{y-S_n}\to0$, and therefore
  $h_y=\sum_{k\ge0}\mathbb E\psi_{y-S_k}$.  The middle region of the three-way
  split is controlled by $\mu^{*n}(I)\to0$ on bounded intervals, which here
  comes from local finiteness of $U$ rather than from almost-sure divergence of
  $S_n$.
\end{lemma}
\begin{proof}\leanok\end{proof}

\begin{lemma}[Key renewal limit for Gaussian increments]
  \label{lem:nd-gaussian-renewal}
  \lean{AbsorptionCutoff.Renewal.tendsto_of_renewalEquation_of_driNorm,
    AbsorptionCutoff.Renewal.tendsto_atTop_zero_of_driNorm,
    AbsorptionCutoff.Renewal.tendsto_angular_of_renewalEquation_of_driNorm,
    AbsorptionCutoff.tendsto_of_renewalEquation_tiltedIncrementLaw,
    AbsorptionCutoff.tendsto_angular_of_renewalEquation_tiltedIncrementLaw}
  \leanok
  \uses{thm:nd-key-renewal-tilted, lem:nd-renewal-terminal}
  Let $\widehat P_{A,N}$ be the tilted kernel, and assume that the forcing has
  finite directly-Riemann-integrable norm. Then
  \[
    \mathcal H_y(\varphi)\longrightarrow
    \frac{1}{\widehat m_{A,N}}
    \Bigl(\int_{\mathbb R}\Psi_t(\mathbf 1)\,\mathrm dt\Bigr)
    \int_{\mathbb S^{N-1}}\varphi\,\mathrm d\bar\sigma_N ,
    \qquad y\to\infty .
  \]
  Angular reconstruction needs no measure on the sphere: the product form of the
  tilted kernel factors the angular integral out as the constant
  $\int\varphi\,\mathrm d\bar\sigma_N$, leaving the same scalar convolution for
  every $\varphi$.
\end{lemma}
\begin{proof}\leanok\end{proof}

\begin{lemma}[Exponential moments below the Cram\'er exponent]
  \label{lem:nd-subcritical-exponential-moments}
  \lean{AbsorptionCutoff.exists_invariant_Kchain_integrable_neg_rpow,
    AbsorptionCutoff.integrable_neg_rpow_of_invariant_Kchain,
    AbsorptionCutoff.integrable_neg_rpow_gaussianEuclideanNorm_of_invariant_Pkernel}
  \leanok
  \uses{lem:nd-gaussian-cramer-exponent}
  Negative moments of the invariant law are finite strictly below the Cram\'er
  exponent, by a Foster-type drift estimate on the squared-radius chain.
\end{lemma}
\begin{proof}\leanok\end{proof}

\begin{lemma}[Nonlinear polar perturbation estimate]
  \label{lem:nd-gaussian-polar-perturbation}
  \lean{AbsorptionCutoff.exists_gaussianPolarPerturbation_bound}
  \leanok
  \uses{lem:nd-gaussian-cramer-exponent}
  For every $\delta\in(0,2]$, the angular parts of the nonlinear Gaussian
  update and its linearization, restricted at log-radius $t$, differ in total
  variation by at most
  \[
    C r^\delta \min\{1,e^{-Nt}\},\qquad 0<r\le\sqrt N,
  \]
  where the exponentially weighted envelope has finite directly Riemann
  integrable norm.  The proof formalizes the transformed density, its Taylor
  estimate, Gaussian tails, the compact moderate-$r$ branch, and contraction
  under ball restriction and angular projection.
\end{lemma}
\begin{proof}\leanok\end{proof}

\begin{definition}[Nonlinear renewal forcing]
  \label{def:nd-nonlinear-forcing}
  \lean{AbsorptionCutoff.nonlinearForcing}
  \leanok
  For an invariant law $\pi_{A,N}$, the nonlinear renewal forcing
  $\Psi_y^\pi(\varphi)$ is the exponentially tilted difference between the
  nonlinear and linearized log-polar threshold integrals.
\end{definition}

\begin{proposition}[Admissibility of the nonlinear forcing]
  \label{prop:nd-forcing-admissibility}
  \lean{AbsorptionCutoff.driNorm_ofReal_abs_nonlinearForcing_ne_top,
    AbsorptionCutoff.continuous_nonlinearForcing,
    AbsorptionCutoff.nonlinearForcing_one_nonneg,
    AbsorptionCutoff.integral_nonlinearForcing_one,
    AbsorptionCutoff.integral_nonlinearForcing_one_pos}
  \leanok
  \uses{def:nd-nonlinear-forcing,lem:nd-gaussian-polar-perturbation,
    lem:nd-subcritical-exponential-moments}
  Every bounded measurable angular test has finite directly Riemann integrable
  norm and continuous dependence on the log-radius level.  At the constant-one
  test the forcing is nonnegative, its total mass has the paper's Tonelli
  representation, and that mass is strictly positive.
\end{proposition}
\begin{proof}\leanok\end{proof}
